\section{OP-Facing Diagnostic Appendix}
\label{sec:v410r01-op-facing-diagnostic-appendix}

\begin{quote}\small
\noindent\textbf{Status.}
This working layer starts the post-\QSCver{shifted.0.16} OP-facing phase. It does not
rewrite the Companion into Python, does not replace MML/MMA/QAM, and does not
claim that executable diagnostics are proof certificates. The Python files
listed here are diagnostic shadows of selected QAM/MML/MMA/ZFS-style objects.
They are included to help test route shifting, backtranslation, route/support
separation, packed-mask recomputation, guard retention, residual localization,
and proof-state provenance before OP~1/4 and OP~5(B) are attacked in the
formal text.
\end{quote}

\subsection{Diagnostic-code status rule}
\label{subsec:v410r01-diagnostic-code-status-rule}

\begin{definition}[Diagnostic-code status rule]
\label{def:v410r01-diagnostic-code-status-rule}
A diagnostic Python file in this appendix is an executable audit aid. It may
serialize or test selected shadows of QAM, MML, MMA, and ZFS-style records, but
it is not itself a formal definition, not a proof certificate, not a physical
simulation, and not a replacement for the paper's mathematical language.
\end{definition}

\begin{remark}[Why Python appears only here]
\label{rem:v410r01-why-python-only-appendix}
Python appears here only because the current shifted language is rigid enough that
some role separations can be tested mechanically. The MML remains the
mathematical router/intertwiner language in the paper. Python only records
diagnostic encodings of selected finite examples.
\end{remark}

\subsection{Diagnostic files included in the local bundle}
\label{subsec:v410r01-diagnostic-files-included}

\begin{center}
\scriptsize
\setlength{\tabcolsep}{2pt}
\begin{tabular}{@{}p{0.25\textwidth}P{0.62\textwidth}@{}}
\toprule
File & Diagnostic role\\
\midrule
\texttt{qam\_normal\_form.py} & builds a small executable shadow of \(N_4(X)\) with route, guards, and status\\
\texttt{route\_mask\_decoder.py} & separates route, support, and packed mask; verifies recomputation after shift\\
\texttt{mml\_op\_router.py} & records OP-facing MML route surfaces and obstruction kind\\
\texttt{mma\_readout\_residual.py} & records readout residuals and promotion guards\\
\texttt{zfs\_snapshot\_manifest.py} & records proof-state provenance/checksum snapshots\\
\texttt{run\_all\_diagnostics.py} & runs the diagnostic examples and writes a JSON report\\
\bottomrule
\end{tabular}
\end{center}

\subsection{QAM normal-form diagnostic excerpt}
\label{subsec:v410r01-qam-normal-form-diagnostic-excerpt}
\begin{verbatim}
#!/usr/bin/env python3
"""
QAM normal-form diagnostic encoding.

This file is not a proof engine and not an implementation of QAM as a new
foundation. It serializes selected shifted/QAM roles
can be inspected consistently.
"""
from __future__ import annotations

from dataclasses import asdict, dataclass
from enum import Enum
import json
from typing import Tuple


class Status(str, Enum):
 OPEN = "open"
 PREPARED = "prepared"
 PROMOTED = "promoted"


class Guard(str, Enum):
 OP_OPEN = "op_open"
 NO_PHYSICAL_CLAIM = "no_physical_claim"
 NEEDS_SELECTOR = "needs_selector"
 NEEDS_BOUNDARY_PROMOTION = "needs_boundary_promotion"


def sigma_4(k: int) -> int:
 """shifted Hilbert shift on inhabited pre-shift tags."""
 if k < 0:
 raise ValueError("invalid pre-shift tag")
 return k + 1


def back_4(a: int) -> int:
 """Backtranslation from inhabited shifted tags."""
 if a < 1:
 raise ValueError("0_shift has no backtranslation")
 return a - 1

\end{verbatim}

\subsection{Route/mask diagnostic excerpt}
\label{subsec:v410r01-route-mask-diagnostic-excerpt}
\begin{verbatim}
#!/usr/bin/env python3
"""
Route/support/packed-mask diagnostic checks.

The key rule is that a packed mask is recomputed from shifted support. It is
not obtained by adding 1 to an integer mask.
"""
from __future__ import annotations

from dataclasses import dataclass, asdict
import json
from typing import FrozenSet, Tuple


def sigma_4(k: int) -> int:
 if k < 0:
 raise ValueError("invalid pre-shift tag")
 return k + 1


def back_4(a: int) -> int:
 if a < 1:
 raise ValueError("0_shift has no backtranslation")
 return a - 1


def packed_mask_from_support(support: FrozenSet[int]) -> int:
 return sum(2 ** k for k in support)


@dataclass(frozen=True)
class RouteMaskRecord:
 pre_route: Tuple[int, ...]
 shifted_route: Tuple[int, ...]
 support: FrozenSet[int]
 shifted_support: FrozenSet[int]
\end{verbatim}

\subsection{MML and MMA diagnostic excerpts}
\label{subsec:v410r01-mml-mma-diagnostic-excerpts}
\begin{verbatim}
#!/usr/bin/env python3
"""
MML_OP diagnostic router.

This is an executable shadow of an MML OP-facing route record. It does not
define MML and does not prove an OP.
"""
from __future__ import annotations

from dataclasses import asdict, dataclass
from enum import Enum
import json
from typing import Tuple


class Branch(str, Enum):
 HORIZON = "horizon"
 BOUNDARY = "boundary"


class ObstructionKind(str, Enum):
 MISSING_SELECTOR = "missing_selector"
 MISSING_BOUNDARY_PROMOTION = "missing_boundary_promotion"
 UNKNOWN = "unknown"


@dataclass(frozen=True)
class MMLOPRoute:
 name: str
 branch: Branch
 pre_route: Tuple[int, ...]
 shifted_route: Tuple[int, ...]
\end{verbatim}

\begin{verbatim}
#!/usr/bin/env python3
"""
MMA_OP readout-residual diagnostic.

This file gives a small data shape for readout residuals. It does not compute
physical predictions; it records which residual/guard still blocks promotion.
"""
from __future__ import annotations

from dataclasses import asdict, dataclass
from enum import Enum
import json
from typing import Optional


class PromotionStatus(str, Enum):
 BLOCKED = "blocked"
 READY_FOR_REVIEW = "ready_for_review"
 PROMOTED = "promoted"


@dataclass(frozen=True)
class MMAReadoutResidual:
 name: str
 readout_target: str
 hc_boundary_condition: str
 comparison_selector: Optional[str]
 residual_surface: str
 promotion_guard: str
 status: PromotionStatus
\end{verbatim}

\subsection{ZFS-style provenance diagnostic excerpt}
\label{subsec:v410r01-zfs-provenance-diagnostic-excerpt}
\begin{verbatim}
#!/usr/bin/env python3
"""
ZFS-style proof-state snapshot manifest.

This is a provenance/audit helper. It is not a mathematical object and not a
proof operator.
"""
from __future__ import annotations

from dataclasses import asdict, dataclass
import hashlib
import json
from pathlib import Path
from typing import Dict, Iterable


@dataclass(frozen=True)
class SnapshotManifest:
 base_version: str
 branch_name: str
 files: Dict[str, str]
 purpose: str
 non_claims: tuple[str, ...]

 def to_json_dict(self) -> dict:
 return asdict(self)


def sha256_file(path: Path) -> str:
 return hashlib.sha256(path.read_bytes()).hexdigest()
\end{verbatim}

\subsection{First diagnostic run record}
\label{subsec:v410r01-first-diagnostic-run-record}

The local diagnostic runner completed successfully. The generated JSON report
is stored in the local bundle as
\[
 \texttt{python\_diagnostics/diagnostic\_run\_report.json}.
\]
The successful run only certifies that the diagnostic examples execute and that
their internal consistency checks pass. It does not close OP~1/4 or OP~5(B).

\begin{theorem}[Diagnostic appendix conservativity]
\label{thm:v410r01-diagnostic-appendix-conservativity}
The local proof line diagnostic-code appendix is conservative over \QSCver{shifted.0.16}
release-final.
\end{theorem}

\begin{proof}
The appendix adds executable diagnostic shadows and records their status as
audit aids. It does not modify \(\Sigma_4\), \(\mathrm B_4\),
\(\mathfrak D_4\), \(\mathcal N_4\), tensor grammar, readout compatibility,
Historical Witness content, or OP-status guards. Therefore it is conservative.
\end{proof}

\begin{remark}[non-claim]
\label{rem:v410r01-non-claim}
Version~4.1.0-does not add new physics, does not close OP~1/4 or OP~5(B),
does not activate primitive mixed tensors, does not claim Turing-completeness,
and does not replace standard QFT/QED/Lifshitz or general-relativistic
calculations.
\end{remark}



\clearpage
\section{First OP Obstruction Extraction}
\label{sec:v410r02-first-op-obstruction-extraction}

\begin{quote}\small
\noindent\textbf{Status.}
This working layer uses the diagnostic scripts from the local proof line against the two
current OP-facing targets OP~1/4 and OP~5(B). The output is an obstruction
extraction report, not a proof closure. It identifies the current named block
that prevents promotion of each OP packet.
\end{quote}

\subsection{Obstruction extraction rule}
\label{subsec:v410r02-obstruction-extraction-rule}

\begin{definition}[Diagnostic obstruction extraction]
\label{def:v410r02-diagnostic-obstruction-extraction}
A diagnostic obstruction extraction sends an OP-facing packet to a finite audit
record containing its branch, pre-shift route, shifted route, backtranslation,
support, shifted support, packed masks, witness link, active normal form,
readout target, missing item, guard, and current status. The record is a
machine-checkable audit surface and not a theorematic closure certificate.
\end{definition}

\subsection{First extracted OP obstruction table}
\label{subsec:v410r02-first-extracted-op-obstruction-table}

\begin{center}
\small
\setlength{\tabcolsep}{3pt}
\begin{tabular}{@{}p{0.16\textwidth}P{0.12\textwidth}P{0.18\textwidth}P{0.23\textwidth}P{0.21\textwidth}@{}}
\toprule
OP packet & branch & active form & obstruction & missing item\\
\midrule
OP~1/4 & horizon & $N_4(\mathrm{OP}_{1/4})$ & missing selector & comparison tensor $C_{1/4}^{\mathrm{sel}}$ \\
OP~5(B) & boundary & $N_4(O\Pi_{5}^{(B)})$ & missing boundary promotion & theorem $\Pi_{5}^{(B)}$ \\
\bottomrule
\end{tabular}
\end{center}

The generated files are included in the local bundle under
\texttt{python\_diagnostics/}: \texttt{op\_obstruction\_report.json},
\texttt{op\_obstruction\_table.md}, and \texttt{op\_obstruction\_table.csv}.

\subsection{Interpretation}
\label{subsec:v410r02-first-op-obstruction-interpretation}

\begin{proposition}[First OP obstruction localization]
\label{prop:v410r02-first-op-obstruction-localization}
Under the current diagnostic encoding, OP~1/4 is reduced to a missing
selector/comparison tensor on the horizon-side route-normalized packet, while
OP~5(B) is reduced to a missing boundary-promotion theorem preserving the
boundary-side stop guard.
\end{proposition}

\begin{proof}
The diagnostic extractor verifies route shifting and backtranslation for each
packet, keeps the Historical Witness link separate from the active normal
form, and refuses to promote a packet when its declared missing item remains
active. The OP~1/4 packet is horizon-side and therefore fails at the declared
selector/comparison requirement. The OP~5(B) packet is boundary-side and
therefore fails at the declared boundary-promotion requirement. These failures
are named obstructions rather than closures.
\end{proof}

\begin{remark}[What the extractor has and has not shown]
\label{rem:v410r02-what-extractor-shows}
The extraction has shown where the current OP-facing packets are blocked in the
shifted diagnostic grammar. It has not proven that the selector/comparison tensor
exists, has not proven the boundary-promotion theorem, and has not promoted any
physical readout.
\end{remark}

\begin{theorem}[First obstruction extraction conservativity]
\label{thm:v410r02-first-obstruction-extraction-conservativity}
The local proof line first OP obstruction extraction is conservative over the local proof line
and over \QSCver{shifted.0.16} release-final.
\end{theorem}

\begin{proof}
The working layer adds a diagnostic extraction report and a reader-facing
summary table. It does not modify \(\Sigma_4\), \(\mathrm B_4\),
\(\mathfrak D_4\), \(\mathcal N_4\), tensor grammar, readout compatibility,
Historical Witness content, or OP-status guards. OP~1/4 and OP~5(B) remain
open. Therefore the layer is conservative.
\end{proof}

\begin{remark}[non-claim]
\label{rem:v410r02-non-claim}
Version~4.1.0-does not close OP~1/4 or OP~5(B). It only names their first
current diagnostic obstruction surfaces.
\end{remark}



\clearpage
\section{Obstruction-to-Task Reduction Layer}
\label{sec:v410r03-obstruction-to-task-reduction-layer}

\begin{quote}\small
\noindent\textbf{Status.}
This working layer consumes the diagnostic obstruction extraction and turns
its two named obstructions into a finite OP-facing task queue. It does not
close OP~1/4 or OP~5(B), does not promote a physical readout, and does not
replace the formal QAM/MML/MMA language by Python. The Python task queue is an
audit and planning surface only.
\end{quote}

\subsection{From obstruction labels to task packets}
\label{subsec:v410r03-from-obstruction-labels-to-task-packets}

\begin{definition}[OP obstruction-to-task packet]
\label{def:v410r03-op-obstruction-to-task-packet}
An OP obstruction-to-task packet is a finite record
\[
 \mathcal T_{\mathrm{OP}}
 =
 (O,B,N_4(O),\tau,T,I,C,G,S),
\]
where \(O\) is the OP object, \(B\) is the branch, \(N_4(O)\) is the active
shifted-namespace normal form, \(\tau\) is the named obstruction, \(T\) is the next task
kind, \(I\) is the required input, \(C\) is the success criterion, \(G\) is the
guard to preserve, and \(S\) is the task status.
\end{definition}

\begin{remark}[Task packets are not proof packets]
\label{rem:v410r03-task-packets-not-proof-packets}
The task packet names the next controlled move. It does not certify that the
move can be completed, and it does not close the corresponding OP.
\end{remark}

\subsection{Extracted task queue}
\label{subsec:v410r03-extracted-task-queue}

The diagnostic task planner reads the file
\[
 \texttt{python\_diagnostics/op\_obstruction\_report.json}
\]
and produces
\[
 \texttt{op\_task\_queue.json},\qquad
 \texttt{op\_task\_queue.md},\qquad
 \texttt{op\_task\_queue.csv}.
\]
The resulting task queue is:

\begin{center}
\small
\setlength{\tabcolsep}{3pt}
\begin{tabular}{@{}p{0.12\textwidth}P{0.12\textwidth}P{0.26\textwidth}P{0.25\textwidth}P{0.12\textwidth}@{}}
\toprule
OP & Task & Kind & Target object & Status\\
\midrule
OP~1/4 & T1 & selector domain/codomain & $C_{1/4}^{\mathrm{sel}}$ & ready \\
OP~1/4 & T2 & backtranslation check & $B_4$-compatible selector & open \\
OP~1/4 & T3 & residual reduction & $R_{1/4}^{\mathrm{post}}$ & open \\
OP~5(B) & T1 & promotion theorem & $\Pi_{5}^{(B)}$ & ready \\
OP~5(B) & T2 & HC/MM compatibility & boundary compatibility check & open \\
OP~5(B) & T3 & promotion residual & $R_{5}^{(B),\mathrm{post}}$ & open \\
\bottomrule
\end{tabular}
\end{center}

\subsection{OP~1/4 reduction}
\label{subsec:v410r03-op14-reduction}

\begin{proposition}[OP~1/4 selector-task reduction]
\label{prop:v410r03-op14-selector-task-reduction}
Given the obstruction record, the first controlled OP~1/4 move is not a
new physical readout and not an OP closure. It is the construction of a
selector/comparison tensor \(C_{1/4}^{\mathrm{sel}}\) with explicit domain,
codomain, route role, and backtranslation compatibility.
\end{proposition}

\begin{proof}
The obstruction is \texttt{missing\_selector} and the missing item is
\texttt{selector\_or\_comparison\_tensor}. The planner therefore
extracts the first task \texttt{OP14-T1}, whose target is the selector or
comparison tensor. Since the associated guard remains \texttt{OP\_OPEN} and
\texttt{no physical claim}, the task is a controlled reduction rather than an
OP closure.
\end{proof}

\subsection{OP~5(B) reduction}
\label{subsec:v410r03-op5b-reduction}

\begin{proposition}[OP~5(B) boundary-promotion task reduction]
\label{prop:v410r03-op5b-boundary-promotion-task-reduction}
Given the obstruction record, the first controlled OP~5(B) move is the
statement of a boundary-promotion theorem \(\Pi_{5}^{(B)}\) preserving the earliest
boundary stop. It is not a physical promotion and not an OP closure.
\end{proposition}

\begin{proof}
The obstruction is \texttt{missing\_boundary\_promotion} and the missing
item is \texttt{boundary\_promotion\_theorem}. The planner therefore
extracts the first task \texttt{OP5B-T1}, whose target is a boundary
promotion theorem. Since the guard requires the stop-preserving boundary
condition and no physical claim, the task cannot be read as a promoted physical
readout or as OP closure.
\end{proof}

\subsection{Next formal move}
\label{subsec:v410r03-next-formal-move}

The next mathematical step after is therefore sharply constrained:
\[
 \text{OP~1/4: construct/test } C_{1/4}^{\mathrm{sel}},
 \qquad
 \text{OP~5(B): state/test } \Pi_{5}^{(B)}.
\]
If \(C_{1/4}^{\mathrm{sel}}\) fails, the failure becomes the next named OP~1/4
obstruction. If \(\Pi_{5}^{(B)}\) fails, the failure becomes the next named OP~5(B)
boundary obstruction. In both cases the OP remains open unless the residual is
formally discharged in the Active Core.

\begin{theorem}[Obstruction-to-task reduction conservativity]
\label{thm:v410r03-obstruction-to-task-reduction-conservativity}
The local proof line task-object formalization is conservative over the local proof line.
\end{theorem}

\begin{proof}
The layer reads the diagnostic obstruction records and produces a
finite task queue. It does not change \(\Sigma_4\), \(\mathrm B_4\),
\(\mathfrak D_4\), \(\mathcal N_4\), tensor grammar, readout compatibility,
Historical Witness content, or OP-status guards. Therefore it is conservative
over the local pass.
\end{proof}

\begin{remark}[non-claim]
\label{rem:v410r03-non-claim}
Version~4.1.0-does not add new physics, does not close OP~1/4 or OP~5(B),
does not activate primitive mixed tensors, does not claim Turing-completeness,
and does not replace standard QFT/QED/Lifshitz or general-relativistic
calculations.
\end{remark}



\clearpage
\section{Task-Object Formalization}
\label{sec:v410r04-task-object-formalization}

\begin{quote}\small
\noindent\textbf{Status.}
This layer turns the task queue into two explicit next target objects. It
is still not an OP closure. It identifies the first formal object needed for
OP~1/4 as a selector/comparison tensor $C_{1/4}^{\mathrm{sel}}$ and the first formal object
needed for OP~5(B) as a boundary-promotion theorem $\Pi_{5}^{(B)}$.
\end{quote}

\subsection{Task-to-object reduction rule}
\label{subsec:v410r04-task-to-object-reduction-rule}

\begin{definition}[Task-to-object reduction]
\label{def:v410r04-task-to-object-reduction}
A task-to-object reduction sends a named obstruction task
\[
 T=(OP,\mathrm{obstruction},\mathrm{missing\ item},G,S)
\]
to a target object
\[
 \mathcal T(T)=(X_T,D_T,E_T,\mathcal C_T,R_T^{post},G_T,S_T),
\]
where $X_T$ is the object to be constructed, $D_T$ is its intended domain,
$E_T$ is its comparison or promotion codomain, $\mathcal C_T$ is the list of
required compatibility conditions, $R_T^{post}$ is the residual left after the
construction, $G_T$ is the retained guard, and $S_T$ is the post-reduction
status.
\end{definition}

\begin{remark}[Why this is mathematical progress]
\label{rem:v410r04-why-this-is-progress}
The local pass/diagnostics named where the packets fail. The reduction names
what must be constructed next. This changes an unlocalized OP obstruction into
a concrete target-object problem.
\end{remark}

\subsection{The OP~1/4 target object $C_{1/4}^{\mathrm{sel}}$}
\label{subsec:v410r04-op14-target-object}

\begin{definition}[OP~1/4 selector/comparison target]
\label{def:v410r04-op14-selector-comparison-target}
Let $X_{14}=N_4(\mathrm{OP}_{1/4})$ be the horizon-side OP~1/4 normal-form packet, with
pre-shift witness route $\rho_{1/4}^{\mathrm{pre}}$ and inhabited shifted route $\rho_{1/4}^{\mathrm{shift}}$.
A first OP~1/4 selector/comparison target is a comparison tensor
\[
 C_{1/4}^{\mathrm{sel}}:\quad
 P(\rho_{1/4}^{\mathrm{pre}})\times P(\rho_{1/4}^{\mathrm{shift}})\longrightarrow E_{1/4}^{\mathrm{sel}}
\]
subject to backtranslation compatibility, guard preservation, and residual
monotonicity $R_{1/4}^{\mathrm{post}}\preccurlyeq R_{1/4}^{\mathrm{pre}}$.
\end{definition}

\begin{proposition}[OP~1/4 task-object reduction]
\label{prop:v410r04-op14-task-object-reduction}
Under the obstruction table, OP~1/4 reduces first to the construction of
$C_{1/4}^{\mathrm{sel}}$ satisfying Definition~\ref{def:v410r04-op14-selector-comparison-target}.
If such a $C_{1/4}^{\mathrm{sel}}$ is constructed, the remaining OP~1/4 obstruction is the
post-selector residual $R_{1/4}^{\mathrm{post}}$, not the original missing-selector
obstruction.
\end{proposition}

\begin{proof}
The task queue assigns OP~1/4 the obstruction \emph{missing selector}. A
selector obstruction is discharged only by an object with a declared domain,
comparison codomain, compatibility conditions, and retained guard. Definition
\ref{def:v410r04-op14-selector-comparison-target} supplies exactly this target.
It preserves backtranslation and guard data and leaves an explicit residual.
Therefore the first obstruction is reduced to constructing $C_{1/4}^{\mathrm{sel}}$, after
which only $R_{1/4}^{\mathrm{post}}$ remains to be analyzed.
\end{proof}

\subsection{The OP~5(B) target object $\Pi_{5}^{(B)}$}
\label{subsec:v410r04-op5b-target-object}

\begin{definition}[OP~5(B) boundary-promotion target]
\label{def:v410r04-op5b-boundary-promotion-target}
Let $Y_{5}^{(B)}=N_4(O\Pi_{5}^{(B)})$ be the boundary-side OP~5(B) normal-form packet with
stop-preserving boundary guard $G_{\partial}$. A first OP~5(B)
boundary-promotion target is a theorem schema $\Pi_{5}^{(B)}$ of the form
\[
 \Pi_{5}^{(B)}:\quad
 (Y_{5}^{(B)},G_{\partial},HC/MM)\longrightarrow
 (Y_{5}^{(B),\mathrm{prom}},G_{\partial},R_{5}^{(B),\mathrm{post}}),
\]
subject to stop preservation, HC/MM compatibility, backtranslation stability,
and residual retention.
\end{definition}

\begin{proposition}[OP~5(B) task-object reduction]
\label{prop:v410r04-op5b-task-object-reduction}
Under the obstruction table, OP~5(B) reduces first to the formulation and
proof of $\Pi_{5}^{(B)}$ satisfying Definition~\ref{def:v410r04-op5b-boundary-promotion-target}.
If such a theorem is obtained, the remaining OP~5(B) obstruction is
$R_{5}^{(B),\mathrm{post}}$ under the retained stop-preserving guard.
\end{proposition}

\begin{proof}
The task queue assigns OP~5(B) the obstruction \emph{missing boundary
promotion}. A boundary-promotion obstruction cannot be discharged by a route
renaming alone, since the stop-preserving guard must survive promotion. The
schema $\Pi_{5}^{(B)}$ states the exact data required for a conservative promotion:
retained guard, HC/MM compatibility, backtranslation stability, and explicit
post-promotion residual. Hence OP~5(B) is reduced to proving $\Pi_{5}^{(B)}$, after
which $R_{5}^{(B),\mathrm{post}}$ is the remaining obstruction.
\end{proof}

\begin{theorem}[task-object formalization conservativity]
\label{thm:v410r04-task-object-formalization-conservativity}
The local proof line task-object formalization is conservative over the local proof line.
\end{theorem}

\begin{proof}
The layer does not claim that $C_{1/4}^{\mathrm{sel}}$ or $\Pi_{5}^{(B)}$ has already been
constructed. It only states their required roles and the residuals that remain
after their construction. OP~1/4 and OP~5(B) therefore remain open, while their
first target objects are now named and constrained.
\end{proof}



\clearpage
\section{Notation and Physics-Style Audit for OP-Facing Objects}
\label{sec:v410r05-notation-physics-style-audit}

\begin{quote}\small
\noindent\textbf{Status.}
This layer audits the OP-facing notation introduced in the local proof chain and separates
mathematical display notation from diagnostic/code-safe identifiers. It does
not change the obstruction analysis, does not close OP~1/4 or OP~5(B), and
does not replace QAM/MML/MMA by Python or by information-theoretic syntax.
\end{quote}

\subsection{Display notation versus diagnostic identifiers}
\label{subsec:v410r05-display-notation-versus-diagnostic-identifiers}

\begin{definition}[Display notation and diagnostic identifier]
\label{def:v410r05-display-notation-diagnostic-identifier}
A display notation is the notation used in mathematical and physical statements.
A diagnostic identifier is a machine-safe name used by Python snippets, JSON
reports, CSV tables, or local audit files. Diagnostic identifiers may be short
and ASCII-safe; display notation must remain mathematically and physically
readable.
\end{definition}

\begin{center}
\small
\setlength{\tabcolsep}{3pt}
\begin{tabular}{@{}p{0.25\textwidth}P{0.28\textwidth}P{0.35\textwidth}@{}}
\toprule
Diagnostic identifier & Mathematical display notation & Role\\
\midrule
\texttt{OP\_1\_4} & \(\mathrm{OP}_{1/4}\) & OP~1/4 normal-form target\\
\texttt{OP\_5B} & \(\mathrm{OP}_{5}^{(B)}\) & OP~5(B) normal-form target\\
\texttt{C\_OP14} & \(C_{1/4}^{\mathrm{sel}}\) & selector/comparison object for OP~1/4\\
\texttt{P\_5B} & \(\Pi_{5}^{(B)}\) & boundary-promotion schema for OP~5(B)\\
\texttt{R\_OP14\_post} & \(R_{1/4}^{\mathrm{post}}\) & post-selector residual\\
\texttt{R\_5B\_post} & \(R_{5}^{(B),\mathrm{post}}\) & post-promotion residual\\
\bottomrule
\end{tabular}
\end{center}

\subsection{Physics-style rule}
\label{subsec:v410r05-physics-style-rule}

\begin{definition}[Physics-style rule for OP-facing notation]
\label{def:v410r05-physics-style-rule}
Information-theoretic and code-facing tools may be used as diagnostic encodings
only after the governing mathematical object has been written in display
notation. In particular, the active text should prefer
\[
 C_{1/4}^{\mathrm{sel}},\qquad
 \Pi_{5}^{(B)},\qquad
 R_{1/4}^{\mathrm{post}},\qquad
 R_{5}^{(B),\mathrm{post}}
\]
over code-like symbols such as \texttt{C\_OP14}, \texttt{P\_5B},
\texttt{R\_OP14}, or \texttt{R\_5B}.
\end{definition}

\begin{remark}[Information theory is a diagnostic layer]
\label{rem:v410r05-information-theory-diagnostic-layer}
The route/mask/identifier layer may look information-theoretic, but the
governing objects remain mathematical and physical readout objects. The
information-theoretic encoding is therefore a compression and audit surface, not
a replacement ontology.
\end{remark}

\subsection{Normalized target-object statements}
\label{subsec:v410r05-normalized-target-object-statements}

With notation normalized, the two target objects read as follows:
\[
 C_{1/4}^{\mathrm{sel}}:
 P(\rho_{1/4}^{\mathrm{pre}})\times P(\rho_{1/4}^{\mathrm{shift}})
 \longrightarrow E_{1/4}^{\mathrm{sel}},
\]
and
\[
 \Pi_{5}^{(B)}:
 (Y_{5}^{(B)},G_{\partial},HC/MM)
 \longrightarrow
 (Y_{5}^{(B),\mathrm{prom}},G_{\partial},R_{5}^{(B),\mathrm{post}}).
\]
Thus OP~1/4 is still reduced to a selector/comparison construction, and
OP~5(B) is still reduced to a boundary-promotion construction, but the display
notation no longer borrows code-style names from the diagnostic layer.

\begin{theorem}[Notation-normalization conservativity]
\label{thm:v410r05-notation-normalization-conservativity}
The notation normalization is conservative over the task-object
formalization.
\end{theorem}

\begin{proof}
The normalization replaces code-like display labels by mathematical display
notation and records a diagnostic-identifier map. It does not change the
underlying route, guard, obstruction, residual, or OP-status records. Therefore
the reductions to a selector/comparison object for OP~1/4 and to a
boundary-promotion schema for OP~5(B) are preserved.
\end{proof}

\begin{remark}[non-claim]
\label{rem:v410r05-non-claim}
Version~4.1.0-does not close OP~1/4 or OP~5(B), does not activate
primitive mixed tensors, does not claim Turing-completeness, and does not
replace standard QFT/QED/Lifshitz or general-relativistic calculations.
\end{remark}



\clearpage
\section{First Selector and Boundary-Promotion Ansatz}
\label{sec:v410r06-first-selector-boundary-promotion-ansatz}

\begin{quote}\small
\noindent\textbf{Status.}
This layer starts the first genuine OP-facing mathematical attack after the
notation/physics-style audit. It does not close OP~1/4 or OP~5(B). It turns
the two normalized target objects,
\[
 C_{1/4}^{\mathrm{sel}}
 \qquad\text{and}\qquad
 \Pi_{5}^{(B)},
\]
into first construction ansaetze: a comparison/fibre-product ansatz for
OP~1/4 and a guard-preserving boundary-promotion ansatz for OP~5(B).
\end{quote}

\subsection{OP~1/4: comparison surface and selector ansatz}
\label{subsec:v410r06-op14-comparison-surface-selector-ansatz}

\begin{definition}[OP~1/4 comparison surface]
\label{def:v410r06-op14-comparison-surface}
Let
\[
 P(\rho_{1/4}^{\mathrm{pre}})
 \qquad\text{and}\qquad
 P(\rho_{1/4}^{\mathrm{shift}})
\]
be the witness-side and shifted-normalized packet surfaces attached to
\(\mathrm{OP}_{1/4}\). A comparison surface for OP~1/4 is an object
\(E_{1/4}^{\mathrm{sel}}\) equipped with two comparison maps
\[
 e_{\mathrm{pre}}:
 P(\rho_{1/4}^{\mathrm{pre}})\longrightarrow E_{1/4}^{\mathrm{sel}},
 \qquad
 e_{\mathrm{shift}}:
 P(\rho_{1/4}^{\mathrm{shift}})\longrightarrow E_{1/4}^{\mathrm{sel}}.
\]
The maps must be read as current shifted diagnostic/readout maps, not as a claim
that the historical witness and active shifted surface are already identical.
\end{definition}

\begin{definition}[Selector/comparison object for OP~1/4]
\label{def:v410r06-selector-comparison-object-op14}
Given a comparison surface \(E_{1/4}^{\mathrm{sel}}\), define the first
selector/comparison ansatz by the fibre product
\[
 C_{1/4}^{\mathrm{sel}}
 :=
 P(\rho_{1/4}^{\mathrm{pre}})
 \times_{E_{1/4}^{\mathrm{sel}}}
 P(\rho_{1/4}^{\mathrm{shift}}).
\]
Equivalently,
\[
 C_{1/4}^{\mathrm{sel}}
 =
 \left\{(x_{\mathrm{pre}},x_{\mathrm{shift}})
 \;\middle|\;
 e_{\mathrm{pre}}(x_{\mathrm{pre}})
 =
 e_{\mathrm{shift}}(x_{\mathrm{shift}})
 \right\}.
\]
Thus \(C_{1/4}^{\mathrm{sel}}\) contains exactly the witness/shifted pairs that
become equal after being read in the same comparison surface.
\end{definition}

\begin{definition}[Post-selector residual]
\label{def:v410r06-post-selector-residual}
The post-selector residual for OP~1/4 is the remaining comparison obstruction
after the selector object has been formed. In the weak set-difference reading,
\[
 R_{1/4}^{\mathrm{post}}
 :=
 E_{1/4}^{\mathrm{sel}}
 \setminus
 \operatorname{im}\bigl(C_{1/4}^{\mathrm{sel}}\to E_{1/4}^{\mathrm{sel}}\bigr).
\]
In a category or additive setting, the same role may be represented by a
cokernel/cocone residual. The present layer requires only the weak residual
role, not a fixed ambient abelian category.
\end{definition}

\begin{lemma}[Selector ansatz reduces the OP~1/4 obstruction]
\label{lem:v410r06-selector-ansatz-reduces-op14-obstruction}
If the comparison surface exists and the fibre product
\(C_{1/4}^{\mathrm{sel}}\) is well-defined under the inherited guards, then the
obstruction \emph{missing selector} is reduced to the post-selector
residual \(R_{1/4}^{\mathrm{post}}\).
\end{lemma}

\begin{proof}
The obstruction identified the absence of a comparison/selector object.
Definitions~\ref{def:v410r06-op14-comparison-surface} and
\ref{def:v410r06-selector-comparison-object-op14} provide such an object as a
fibre product over a common comparison surface. Therefore the original absence
of a selector is discharged at the level of ansatz formation. What remains is
not the same obstruction but the part of the comparison surface not reached by
that selector, namely \(R_{1/4}^{\mathrm{post}}\).
\end{proof}

\begin{remark}[Why this is not yet OP~1/4 closure]
\label{rem:v410r06-selector-not-op14-closure}
The selector ansatz does not prove that \(R_{1/4}^{\mathrm{post}}=\varnothing\)
and does not prove that all OP~1/4 guards are discharged. It only turns the
vague obstruction ``missing selector'' into a precise residual target.
\end{remark}

\subsection{OP~5(B): guard-preserving boundary-promotion ansatz}
\label{subsec:v410r06-op5b-boundary-promotion-ansatz}

\begin{definition}[Boundary-promotion ansatz for OP~5(B)]
\label{def:v410r06-boundary-promotion-ansatz-op5b}
Let \(Y_{5}^{(B)}\) be the boundary-side object associated with
\(\mathrm{OP}_{5}^{(B)}\), let \(G_{\partial}\) be the inherited
stop-preserving boundary guard, and let \(\mathcal R_{HC/MM}\) denote the
HC/MM readout interface. A first boundary-promotion ansatz is an operation
\[
 \Pi_{5}^{(B)}:
 (Y_{5}^{(B)},G_{\partial},HC/MM)
 \longrightarrow
 (Y_{5}^{(B),\mathrm{prom}},G_{\partial},R_{5}^{(B),\mathrm{post}})
\]
satisfying the guard-preservation condition
\[
 G_{\partial}\circ \Pi_{5}^{(B)}=G_{\partial}
\]
where both sides are read on the admissible boundary-side domain.
\end{definition}

\begin{definition}[HC/MM compatibility defect]
\label{def:v410r06-hc-mm-compatibility-defect}
The HC/MM compatibility defect of \(\Pi_{5}^{(B)}\) is the obstruction to the
commutation relation
\[
 \mathcal R_{HC/MM}\circ \Pi_{5}^{(B)}
 \sim_{\partial}
 \Pi_{5}^{(B)}\circ \mathcal R_{HC/MM},
\]
where \(\sim_{\partial}\) denotes boundary-guard-compatible comparison. The
post-promotion residual is written
\[
 R_{5}^{(B),\mathrm{post}}
 :=
 \operatorname{Obstr}\left(
 G_{\partial}\circ \Pi_{5}^{(B)}=G_{\partial},
 \;
 \mathcal R_{HC/MM}\circ \Pi_{5}^{(B)}
 \sim_{\partial}
 \Pi_{5}^{(B)}\circ \mathcal R_{HC/MM}
 \right).
\]
\end{definition}

\begin{lemma}[Boundary-promotion ansatz reduces the OP~5(B) obstruction]
\label{lem:v410r06-boundary-promotion-ansatz-reduces-op5b-obstruction}
If \(\Pi_{5}^{(B)}\) is defined on the admissible boundary-side domain and
preserves \(G_{\partial}\), then the obstruction \emph{missing boundary
promotion} is reduced to the post-promotion residual
\(R_{5}^{(B),\mathrm{post}}\).
\end{lemma}

\begin{proof}
The obstruction named the absence of a boundary-promotion mechanism. A
map \(\Pi_{5}^{(B)}\) satisfying Definition~\ref{def:v410r06-boundary-promotion-ansatz-op5b}
provides such a mechanism while preserving the inherited boundary guard. The
remaining obstruction is therefore not the absence of promotion itself, but the
failure, if any, of guard preservation and HC/MM compatibility. This remaining
obstruction is precisely \(R_{5}^{(B),\mathrm{post}}\).
\end{proof}

\begin{remark}[Physical reading]
\label{rem:v410r06-op5b-physical-reading}
The OP~5(B) ansatz keeps the physical boundary reading visible. It does not
promote a boundary object by erasing the boundary stop; it asks whether a
promotion through the HC/MM readout interface can be made while the
stop-preserving guard survives.
\end{remark}

\subsection{First combined reduction statement}
\label{subsec:v410r06-first-combined-reduction-statement}

\begin{theorem}[First OP-facing obstruction reduction]
\label{thm:v410r06-first-op-facing-obstruction-reduction}
In the ansatz layer, the two the local pass/obstruction labels reduce as follows:
\[
\begin{array}{rcl}
\mathrm{OP}_{1/4}:\quad
\text{missing selector}
&\leadsto&
R_{1/4}^{\mathrm{post}},\\[0.3em]
\mathrm{OP}_{5}^{(B)}:\quad
\text{missing boundary promotion}
&\leadsto&
R_{5}^{(B),\mathrm{post}}.
\end{array}
\]
The reduction is conditional on the formation of
\(C_{1/4}^{\mathrm{sel}}\) and \(\Pi_{5}^{(B)}\), respectively.
\end{theorem}

\begin{proof}
The first line follows from Lemma~\ref{lem:v410r06-selector-ansatz-reduces-op14-obstruction}.
The second line follows from Lemma~\ref{lem:v410r06-boundary-promotion-ansatz-reduces-op5b-obstruction}.
Both reductions are conditional because neither residual is proven to vanish in
this layer.
\end{proof}

\begin{remark}[non-claim]
\label{rem:v410r06-non-claim}
Version~4.1.0-does not close OP~1/4 or OP~5(B). It replaces the first
named obstructions by sharper residual targets. It also does not activate
primitive mixed tensors, does not claim Turing-completeness, and does not
replace standard QFT/QED/Lifshitz or general-relativistic calculations.
\end{remark}



\clearpage
\section{First Selector Existence Reduction for OP~1/4}
\label{sec:v410r07-selector-existence-reduction-op14}

\begin{quote}\small
\noindent\textbf{Status.}
This layer continues the OP-facing phase by attacking the first OP~1/4 target
\(C_{1/4}^{\mathrm{sel}}\). The point is deliberately modest: the selector is
constructed as a fiber product/equalizer object whenever the two comparison
readouts to a common surface exist. This removes the purely formal
``missing-selector'' obstruction, but it does not close OP~1/4. The remaining
questions are non-emptiness, adequacy, guard strength, and post-selector
residual control.
\end{quote}

\subsection{Comparison surface and readout maps}
\label{subsec:v410r07-comparison-surface-readout-maps}

\begin{definition}[OP~1/4 comparison datum]
\label{def:v410r07-op14-comparison-datum}
An OP~1/4 comparison datum is a triple
\[
 \mathfrak E_{1/4}
 :=
 \bigl(E_{1/4}^{\mathrm{sel}},
 e_{\mathrm{pre}},
 e_{\mathrm{shift}}\bigr)
\]
where \(E_{1/4}^{\mathrm{sel}}\) is a comparison surface and
\[
 e_{\mathrm{pre}}:
 P(\rho_{1/4}^{\mathrm{pre}})
 \longrightarrow
 E_{1/4}^{\mathrm{sel}},
 \qquad
 e_{\mathrm{shift}}:
 P(\rho_{1/4}^{\mathrm{shift}})
 \longrightarrow
 E_{1/4}^{\mathrm{sel}}
\]
are the pre-shift/witness-side and shifted-normalized readout maps.
\end{definition}

\begin{definition}[Selector as fiber product]
\label{def:v410r07-selector-as-fiber-product}
Given a comparison datum \(\mathfrak E_{1/4}\), define
\[
 C_{1/4}^{\mathrm{sel}}
 :=
 P(\rho_{1/4}^{\mathrm{pre}})
 \times_{E_{1/4}^{\mathrm{sel}}}
 P(\rho_{1/4}^{\mathrm{shift}})
\]
as the set of compatible pairs
\[
 C_{1/4}^{\mathrm{sel}}
 =
 \left\{
 (x_{\mathrm{pre}},x_{\mathrm{shift}})
 \ \middle|\
 e_{\mathrm{pre}}(x_{\mathrm{pre}})
 =
 e_{\mathrm{shift}}(x_{\mathrm{shift}})
 \right\}.
\]
Equivalently, it is the equalizer of the two maps
\[
 e_{\mathrm{pre}}\circ \pi_{\mathrm{pre}},
 \quad
 e_{\mathrm{shift}}\circ \pi_{\mathrm{shift}}
 :
 P(\rho_{1/4}^{\mathrm{pre}})
 \times
 P(\rho_{1/4}^{\mathrm{shift}})
 \longrightarrow
 E_{1/4}^{\mathrm{sel}}.
\]
\end{definition}

\begin{lemma}[Formal selector existence]
\label{lem:v410r07-formal-selector-existence}
If \(P(\rho_{1/4}^{\mathrm{pre}})\), \(P(\rho_{1/4}^{\mathrm{shift}})\),
\(E_{1/4}^{\mathrm{sel}}\), \(e_{\mathrm{pre}}\), and \(e_{\mathrm{shift}}\)
exist as ordinary shifted-admissible set/readout objects, then
\(C_{1/4}^{\mathrm{sel}}\) exists as a shifted-admissible selector object.
\end{lemma}

\begin{proof}
The cartesian product
\(P(\rho_{1/4}^{\mathrm{pre}})\times P(\rho_{1/4}^{\mathrm{shift}})\)
exists in the underlying set-compatible universe. The two maps to
\(E_{1/4}^{\mathrm{sel}}\) determine an equality condition. By separation, the
subset of pairs satisfying that equality condition exists. This subset is
exactly \(C_{1/4}^{\mathrm{sel}}\). Since the construction uses only product,
composition, equality, and separation, it is conservative over the shifted-admissible
QAM/ZF-compatible setting.
\end{proof}

\subsection{Backtranslation and guards}
\label{subsec:v410r07-backtranslation-and-guards}

\begin{definition}[Backtranslation-compatible comparison datum]
\label{def:v410r07-backtranslation-compatible-comparison-datum}
A comparison datum \(\mathfrak E_{1/4}\) is backtranslation-compatible if
there is a backtranslation readout \(B_4\) such that the diagram
\[
\begin{array}{ccc}
P(\rho_{1/4}^{\mathrm{shift}}) & \xrightarrow{e_{\mathrm{shift}}}
& E_{1/4}^{\mathrm{sel}}\\
\downarrow B_4 && \downarrow B_E\\
P(\rho_{1/4}^{\mathrm{pre}}) & \xrightarrow{e_{\mathrm{pre}}}
& E_{1/4}^{\mathrm{sel,pre}}
\end{array}
\]
commutes up to the declared comparison tolerance of the OP~1/4 route.
\end{definition}

\begin{lemma}[Guard pullback]
\label{lem:v410r07-guard-pullback}
Let \(G_{1/4}\) be a guard predicate on both packet surfaces and suppose
\(e_{\mathrm{pre}}\) and \(e_{\mathrm{shift}}\) are guard-respecting. Then
\(C_{1/4}^{\mathrm{sel}}\) inherits the pulled-back guard
\[
 G_{1/4}^{\mathrm{sel}}(x_{\mathrm{pre}},x_{\mathrm{shift}})
 :=
 G_{1/4}(x_{\mathrm{pre}})
 \wedge
 G_{1/4}(x_{\mathrm{shift}}).
\]
\end{lemma}

\begin{proof}
Every element of \(C_{1/4}^{\mathrm{sel}}\) is a pair from the two guarded
packet surfaces. Defining the selector guard by conjunction of the two
component guards therefore produces a well-defined predicate on the fiber
product. Because the readout maps are guard-respecting, equality in
\(E_{1/4}^{\mathrm{sel}}\) does not erase the inherited guard data.
\end{proof}

\subsection{Residual reduction}
\label{subsec:v410r07-residual-reduction}

\begin{definition}[Post-selector residual]
\label{def:v410r07-post-selector-residual}
The post-selector residual is
\[
 R_{1/4}^{\mathrm{post}}
 :=
 E_{1/4}^{\mathrm{sel}}
 \setminus
 \operatorname{im}
 \bigl(C_{1/4}^{\mathrm{sel}}\to E_{1/4}^{\mathrm{sel}}\bigr),
\]
or, in a non-set-subtractive readout context, the corresponding obstruction
object measuring the part of \(E_{1/4}^{\mathrm{sel}}\) not covered by
compatible selector pairs.
\end{definition}

\begin{proposition}[Missing selector reduces to adequacy and residual control]
\label{prop:v410r07-missing-selector-reduces-to-adequacy}
Under the comparison-datum hypotheses of
Definition~\ref{def:v410r07-op14-comparison-datum}, the OP~1/4 obstruction
``missing selector'' is replaced by the following narrower tasks:
\[
 \mathrm{NonEmpty}(C_{1/4}^{\mathrm{sel}}),
 \qquad
 \mathrm{Adequate}(C_{1/4}^{\mathrm{sel}}),
 \qquad
 \mathrm{Control}(R_{1/4}^{\mathrm{post}}).
\]
\end{proposition}

\begin{proof}
Lemma~\ref{lem:v410r07-formal-selector-existence} supplies the formal selector
object as soon as the comparison datum exists. Therefore the obstruction is no
longer the absence of a definable selector. What remains is whether the
selector has enough elements or enough structural coverage to compare the two
packet surfaces, and whether the part not covered by the selector is controlled
by \(R_{1/4}^{\mathrm{post}}\). These are precisely non-emptiness, adequacy,
and residual-control tasks.
\end{proof}

\begin{remark}[OP~1/4 status after]
\label{rem:v410r07-op14-status}
This does not close OP~1/4. It does, however, make the formal obstruction
smaller. The next OP~1/4 step is not ``find any selector'', but rather prove
that the fiber-product selector is nonempty or adequate for the OP~1/4 route and
then reduce \(R_{1/4}^{\mathrm{post}}\).
\end{remark}

\begin{theorem}[conservativity]
\label{thm:v410r07-conservativity}
The local proof line selector-existence reduction is conservative over the prior
\QSCver{shifted.1.0} working line.
\end{theorem}

\begin{proof}
The layer constructs \(C_{1/4}^{\mathrm{sel}}\) by ordinary fiber-product /
equalizer logic under stated hypotheses. It does not assert non-emptiness,
adequacy, OP closure, physical promotion, or any replacement of standard
physical calculations. Therefore it is conservative.
\end{proof}



\clearpage
\section{Selector Non-Emptiness and Anchor Criterion for OP~1/4}
\label{sec:v410r08-selector-nonempty-anchor-criterion}

\begin{quote}\small
\noindent\textbf{Status.}
The layer constructed the selector \(C_{1/4}^{\mathrm{sel}}\) as a formal
fiber product once a comparison datum exists. The present layer addresses
the next obstruction: whether that selector is nonempty, and when it is
structurally adequate. The result is still not an OP~1/4 closure. It reduces
non-emptiness to an anchor/intersection condition inside the comparison surface
\(E_{1/4}^{\mathrm{sel}}\).
\end{quote}

\subsection{Anchor set}
\label{subsec:v410r08-anchor-set}

\begin{definition}[OP~1/4 selector anchor set]
\label{def:v410r08-op14-selector-anchor-set}
Given the comparison datum
\[
 \mathfrak E_{1/4}
 =
 (E_{1/4}^{\mathrm{sel}},e_{\mathrm{pre}},e_{\mathrm{shift}}),
\]
define the selector anchor set by
\[
 A_{1/4}^{\mathrm{sel}}
 :=
 \operatorname{im}(e_{\mathrm{pre}})
 \cap
 \operatorname{im}(e_{\mathrm{shift}})
 \subseteq
 E_{1/4}^{\mathrm{sel}}.
\]
An element \(a\in A_{1/4}^{\mathrm{sel}}\) is a common comparison anchor if it
is reached both from the pre-shift/witness-side packet surface and from the
shifted-normalized packet surface.
\end{definition}

\begin{lemma}[Non-emptiness criterion]
\label{lem:v410r08-nonempty-criterion}
The selector \(C_{1/4}^{\mathrm{sel}}\) is nonempty if and only if the anchor
set \(A_{1/4}^{\mathrm{sel}}\) is nonempty.
\end{lemma}

\begin{proof}
If \((x_{\mathrm{pre}},x_{\mathrm{shift}})\in C_{1/4}^{\mathrm{sel}}\), then
\[
 e_{\mathrm{pre}}(x_{\mathrm{pre}})
 =
 e_{\mathrm{shift}}(x_{\mathrm{shift}})=a
\]
for some \(a\in E_{1/4}^{\mathrm{sel}}\). Hence
\(a\in\operatorname{im}(e_{\mathrm{pre}})\cap\operatorname{im}(e_{\mathrm{shift}})\),
so \(A_{1/4}^{\mathrm{sel}}\neq\varnothing\). Conversely, if
\(a\in A_{1/4}^{\mathrm{sel}}\), then there exist
\(x_{\mathrm{pre}}\) and \(x_{\mathrm{shift}}\) with
\(e_{\mathrm{pre}}(x_{\mathrm{pre}})=a=e_{\mathrm{shift}}(x_{\mathrm{shift}})\).
Thus \((x_{\mathrm{pre}},x_{\mathrm{shift}})\in C_{1/4}^{\mathrm{sel}}\).
\end{proof}

\subsection{Adequacy}
\label{subsec:v410r08-adequacy}

\begin{definition}[Anchor adequacy]
\label{def:v410r08-anchor-adequacy}
The anchor set \(A_{1/4}^{\mathrm{sel}}\) is adequate for OP~1/4 if every
comparison-relevant component of \(E_{1/4}^{\mathrm{sel}}\) meets
\(A_{1/4}^{\mathrm{sel}}\). Equivalently, if
\(\pi_0^{\mathrm{rel}}(E_{1/4}^{\mathrm{sel}})\) denotes the finite family of
comparison-relevant components, then
\[
 \forall K\in\pi_0^{\mathrm{rel}}(E_{1/4}^{\mathrm{sel}}):
 \quad
 K\cap A_{1/4}^{\mathrm{sel}}\neq\varnothing.
\]
\end{definition}

\begin{lemma}[Adequacy transfers to selector coverage]
\label{lem:v410r08-adequacy-transfers-to-selector-coverage}
If \(A_{1/4}^{\mathrm{sel}}\) is adequate, then
\(C_{1/4}^{\mathrm{sel}}\) has at least one compatible selector pair over every
comparison-relevant component of \(E_{1/4}^{\mathrm{sel}}\).
\end{lemma}

\begin{proof}
Let \(K\) be a comparison-relevant component. By adequacy, choose
\(a\in K\cap A_{1/4}^{\mathrm{sel}}\). Since \(a\) is in the anchor set, it
has a preimage under both \(e_{\mathrm{pre}}\) and \(e_{\mathrm{shift}}\). The two
preimages form a compatible pair in \(C_{1/4}^{\mathrm{sel}}\) lying over
\(a\in K\).
\end{proof}

\subsection{Residual localization}
\label{subsec:v410r08-residual-localization}

\begin{definition}[Anchor residual]
\label{def:v410r08-anchor-residual}
The anchor residual is
\[
 R_{1/4}^{\mathrm{anc}}
 :=
 E_{1/4}^{\mathrm{sel}}
 \setminus
 A_{1/4}^{\mathrm{sel}},
\]
or, in non-subtractive readout language, the obstruction object measuring the
comparison surface not hit by both readout images.
\end{definition}

\begin{proposition}[Post-selector residual is controlled by anchor residual]
\label{prop:v410r08-post-selector-controlled-by-anchor-residual}
The post-selector residual satisfies the schematic containment
\[
 R_{1/4}^{\mathrm{post}}
 \preceq
 R_{1/4}^{\mathrm{anc}},
\]
where \(\preceq\) denotes the current readout-level obstruction comparison.
In particular, if \(A_{1/4}^{\mathrm{sel}}\) is adequate and the remaining
non-anchor part is readout-irrelevant, then the selector obstruction is
discharged at the comparison level.
\end{proposition}

\begin{proof}
The image of the selector in the comparison surface is exactly the common image
part reached by both readouts, i.e. the anchor set, up to the declared readout
identification. Therefore anything not covered by the selector is contained in
or represented by the complement/non-anchor obstruction. Thus the post-selector
residual is controlled by the anchor residual. Adequacy removes the
comparison-relevant part of this obstruction; any remaining part must be
declared readout-irrelevant before it can be ignored.
\end{proof}

\subsection{Reduced OP~1/4 target after}
\label{subsec:v410r08-reduced-op14-target}

\begin{theorem}[OP~1/4 selector obstruction reduction]
\label{thm:v410r08-op14-selector-obstruction-reduction}
Under the comparison-datum hypotheses, the OP~1/4 selector obstruction reduces
to the following anchor tasks:
\[
 A_{1/4}^{\mathrm{sel}}\neq\varnothing,
 \qquad
 \mathrm{Adequate}(A_{1/4}^{\mathrm{sel}}),
 \qquad
 \mathrm{Control}(R_{1/4}^{\mathrm{anc}}).
\]
\end{theorem}

\begin{proof}
By Lemma~\ref{lem:v410r08-nonempty-criterion}, selector non-emptiness is
equivalent to anchor non-emptiness. By
Lemma~\ref{lem:v410r08-adequacy-transfers-to-selector-coverage}, anchor adequacy
gives selector coverage over the comparison-relevant components. By
Proposition~\ref{prop:v410r08-post-selector-controlled-by-anchor-residual}, the
post-selector residual is controlled by the anchor residual. Hence the
selector obstruction is reduced to anchor non-emptiness, anchor adequacy, and
anchor residual control.
\end{proof}

\begin{remark}[OP~1/4 status after]
\label{rem:v410r08-op14-status-after-r08}
OP~1/4 remains open. However, the obstruction has been sharpened again:
instead of asking whether a selector exists in the abstract, the next step is to
show that the two readout images have a common and adequate anchor set inside
\(E_{1/4}^{\mathrm{sel}}\).
\end{remark}



\clearpage
\section{Anchor Adequacy and Component-Cover Criterion for OP~1/4}
\label{sec:v410r09-anchor-adequacy-component-cover}

\begin{quote}\small
\noindent\textbf{Status.}
The layer reduced selector non-emptiness to the non-emptiness of the anchor
set \(A_{1/4}^{\mathrm{sel}}\). The present layer attacks the next
question: when is the anchor set adequate? The answer is expressed as a
component-cover condition on the comparison-relevant part of
\(E_{1/4}^{\mathrm{sel}}\). This does not close OP~1/4; it converts the
adequacy task into a finite route/component coverage task.
\end{quote}

\subsection{Relevant component system}
\label{subsec:v410r09-relevant-component-system}

\begin{definition}[Comparison-relevant component system]
\label{def:v410r09-comparison-relevant-component-system}
Let \(E_{1/4}^{\mathrm{sel}}\) be the OP~1/4 comparison surface. A
comparison-relevant component system is a finite family
\[
 \mathcal K_{1/4}
 =
 \{K_1,\ldots,K_m\}
\]
of subregions, strata, or readout components of \(E_{1/4}^{\mathrm{sel}}\)
such that every comparison-relevant obstruction in the OP~1/4 selector problem
is represented in at least one \(K_i\).
\end{definition}

\begin{remark}[Why finite components are allowed here]
\label{rem:v410r09-why-finite-components}
The component system is not asserted to exhaust the whole comparison surface.
It only exhausts the comparison-relevant part selected by the current OP~1/4
route and its guards. This keeps the criterion finite and OP-facing.
\end{remark}

\subsection{Anchor cover}
\label{subsec:v410r09-anchor-cover}

\begin{definition}[Anchor cover]
\label{def:v410r09-anchor-cover}
The anchor set \(A_{1/4}^{\mathrm{sel}}\) is a component cover of
\(\mathcal K_{1/4}\) if
\[
 \forall K_i\in\mathcal K_{1/4}:
 \qquad
 K_i\cap A_{1/4}^{\mathrm{sel}}\neq\varnothing.
\]
Equivalently, the cover defect is empty, where
\[
 D_{1/4}^{\mathrm{cov}}
 :=
 \{K_i\in\mathcal K_{1/4}\mid
 K_i\cap A_{1/4}^{\mathrm{sel}}=\varnothing\}.
\]
\end{definition}

\begin{lemma}[Component cover implies adequacy]
\label{lem:v410r09-component-cover-implies-adequacy}
If \(A_{1/4}^{\mathrm{sel}}\) is a component cover of
\(\mathcal K_{1/4}\), then \(A_{1/4}^{\mathrm{sel}}\) is adequate relative to
the comparison-relevant component system.
\end{lemma}

\begin{proof}
Adequacy relative to \(\mathcal K_{1/4}\) means that every comparison-relevant
component contains at least one common anchor. This is exactly the component
cover condition. Hence the condition implies adequacy.
\end{proof}

\subsection{Route-lifted anchors}
\label{subsec:v410r09-route-lifted-anchors}

\begin{definition}[Route-lifted anchor]
\label{def:v410r09-route-lifted-anchor}
A route-lifted anchor over \(K_i\) is a triple
\[
 \alpha_i
 =
 (x_i^{\mathrm{pre}},x_i^{\mathrm{shift}},a_i)
\]
such that
\[
 a_i\in K_i,\qquad
 e_{\mathrm{pre}}(x_i^{\mathrm{pre}})=a_i
 =
 e_{\mathrm{shift}}(x_i^{\mathrm{shift}}),
\]
and the pair \((x_i^{\mathrm{pre}},x_i^{\mathrm{shift}})\) respects the inherited
OP~1/4 route guard.
\end{definition}

\begin{proposition}[Route-lifted anchors certify component cover]
\label{prop:v410r09-route-lifted-anchors-certify-component-cover}
If every \(K_i\in\mathcal K_{1/4}\) admits a route-lifted anchor
\(\alpha_i\), then \(A_{1/4}^{\mathrm{sel}}\) is a component cover of
\(\mathcal K_{1/4}\).
\end{proposition}

\begin{proof}
For each \(K_i\), the anchor \(\alpha_i\) contains a point \(a_i\in K_i\) that
lies in the image of both readout maps. Therefore
\(a_i\in K_i\cap A_{1/4}^{\mathrm{sel}}\), so the intersection is nonempty for
every \(K_i\).
\end{proof}

\subsection{Residual after component cover}
\label{subsec:v410r09-residual-after-component-cover}

\begin{definition}[Cover residual]
\label{def:v410r09-cover-residual}
The cover residual is
\[
 R_{1/4}^{\mathrm{cov}}
 :=
 D_{1/4}^{\mathrm{cov}}
\]
together with any residual guard data attached to components not covered by a
route-lifted anchor.
\end{definition}

\begin{proposition}[Anchor residual reduces to cover residual]
\label{prop:v410r09-anchor-residual-reduces-to-cover-residual}
Relative to the component system \(\mathcal K_{1/4}\), anchor residual control
reduces to cover-residual control:
\[
 \mathrm{Control}(R_{1/4}^{\mathrm{anc}})
 \quad\leadsto\quad
 \mathrm{Control}(R_{1/4}^{\mathrm{cov}}).
\]
If \(D_{1/4}^{\mathrm{cov}}=\varnothing\), then the comparison-relevant anchor
residual is discharged.
\end{proposition}

\begin{proof}
The component system records all comparison-relevant regions. The only
comparison-relevant part not controlled by anchors is represented by components
which do not meet the anchor set. These are precisely the elements of
\(D_{1/4}^{\mathrm{cov}}\). Thus controlling the anchor residual relative to
the selected component system is equivalent to controlling the cover defect.
When the cover defect is empty, every relevant component has an anchor and the
component-level anchor residual vanishes.
\end{proof}

\subsection{Reduced OP~1/4 target after}
\label{subsec:v410r09-reduced-op14-target-after-r09}

\begin{theorem}[OP~1/4 adequacy reduction to finite cover]
\label{thm:v410r09-op14-adequacy-reduction-to-finite-cover}
Under the comparison and anchor hypotheses, OP~1/4 selector adequacy is
reduced to the finite cover condition
\[
 D_{1/4}^{\mathrm{cov}}=\varnothing,
\]
or equivalently to the existence of route-lifted anchors
\[
 \alpha_i=(x_i^{\mathrm{pre}},x_i^{\mathrm{shift}},a_i)
\]
for every comparison-relevant component \(K_i\in\mathcal K_{1/4}\).
\end{theorem}

\begin{proof}
By Lemma~\ref{lem:v410r09-component-cover-implies-adequacy}, a component cover
is exactly adequacy relative to the selected comparison-relevant components.
By Proposition~\ref{prop:v410r09-route-lifted-anchors-certify-component-cover},
route-lifted anchors certify this cover. By
Proposition~\ref{prop:v410r09-anchor-residual-reduces-to-cover-residual}, the
remaining anchor residual is represented by the cover defect. Hence the
adequacy task reduces to the finite cover condition \(D_{1/4}^{\mathrm{cov}}=\varnothing\)
or to the equivalent route-lifted-anchor certificate.
\end{proof}

\begin{remark}[OP~1/4 status after]
\label{rem:v410r09-op14-status-after-r09}
OP~1/4 is still open. However, the obstruction is now much more concrete:
find route-lifted anchors for the finitely selected comparison-relevant
components, or identify which components remain in
\(D_{1/4}^{\mathrm{cov}}\).
\end{remark}



\clearpage
\section{Clock-Anchored Canonical Optimal Route Principle}
\label{sec:v410r10-clock-anchored-canonical-optimal-route}

\begin{quote}\small
\noindent\textbf{Status.}
This layer records the clock-anchored optimal-route principle suggested by the
OP-facing analysis. It changes the problem shape: once the clock anchor and a
target information locus are fixed, the remaining first task is to prove that at
least one admissible route exists. Multiple geometric drawings of the same
canonical route are not treated as different optimal routes; they are identified
by the Stone--von Neumann/Teichmuller/QAM normal-form reading. This does not
close OP~1/4, OP~4, or OP~5. It supplies a stronger route principle for the
next OP reductions.
\end{quote}

\subsection{Clock anchor and admissible route space}
\label{subsec:v410r10-clock-anchor-admissible-route-space}

\begin{definition}[Clock anchor]
\label{def:v410r10-clock-anchor}
The clock anchor is the distinguished known reference information point
\[
 q_{\mathrm{clk}}
\]
against which OP-facing routes are read. It is not an optional coordinate
choice; it is the fixed reference supplied by the taktgeber/clock structure of
the system.
\end{definition}

\begin{definition}[Admissible clock-to-target routes]
\label{def:v410r10-admissible-clock-to-target-routes}
For a target information locus \(q_{\mathrm{tar}}\), let
\[
 \mathcal R_{\mathrm{adm}}(q_{\mathrm{clk}},q_{\mathrm{tar}})
\]
denote the class of router-admissible routes from the clock anchor to the target
which respect the currently declared QAM/MML/MMA tensor interfaces and all
active guards. If a boundary guard \(G_{\partial}\) is present, the admissible
class is restricted to routes satisfying the guard-preservation condition.
\end{definition}

\begin{definition}[Route action and canonical equivalence]
\label{def:v410r10-route-action-canonical-equivalence}
A route action is a ranked route functional
\[
 \mathcal A:
 \mathcal R_{\mathrm{adm}}(q_{\mathrm{clk}},q_{\mathrm{tar}})
 \longrightarrow
 \Lambda
\]
where \(\Lambda\) is the ordered cost/rank set selected by the current
OP-facing readout. Two routes \(\rho_1,\rho_2\) are canonically equivalent,
written
\[
 \rho_1\sim_{\mathrm{can}}\rho_2,
\]
if they represent the same clock-anchored information transport after the
Stone--von Neumann representation identification, the Teichmuller extremal-route
normalization, and the QAM/MML normal-form readout.
\end{definition}

\begin{remark}[Uniqueness means canonical uniqueness]
\label{rem:v410r10-uniqueness-means-canonical-uniqueness}
The principle does not claim that there cannot be two visually different
geometric curves. It claims that, once they realize the same canonical
clock-anchored information transport, they are one route in the mathematical
sense relevant here. Different drawings can therefore represent the same
canonical optimal route.
\end{remark}

\subsection{Canonical optimal route}
\label{subsec:v410r10-canonical-optimal-route}

\begin{definition}[Canonical optimal route class]
\label{def:v410r10-canonical-optimal-route-class}
Assume \(\mathcal R_{\mathrm{adm}}(q_{\mathrm{clk}},q_{\mathrm{tar}})\neq
\varnothing\). A canonical optimal route class is
\[
 [\rho_{\mathrm{opt}}]_{\mathrm{can}}
 :=
 \operatorname*{argmin}_{[\rho]_{\mathrm{can}}
 \subseteq
 \mathcal R_{\mathrm{adm}}(q_{\mathrm{clk}},q_{\mathrm{tar}})}
 \mathcal A([\rho]_{\mathrm{can}}),
\]
where the minimization is taken after canonical equivalence has identified
routes representing the same Stone--von Neumann/Teichmuller/QAM-normalized
information transport.
\end{definition}

\begin{definition}[QAM/MML canonical representative]
\label{def:v410r10-qam-mml-canonical-representative}
The canonical representative
\[
 \rho_{\mathrm{opt}}^{\mathrm{can}}
\]
is the QAM/MML normal-form representative of
\([\rho_{\mathrm{opt}}]_{\mathrm{can}}\). It is the route object used by the
router language, diagnostic comparisons, and OP-facing residual tests.
\end{definition}

\begin{theorem}[Clock-anchored canonical optimal route principle]
\label{thm:v410r10-clock-anchored-canonical-optimal-route-principle}
Suppose the admissible clock-to-target route space is nonempty and finite,
discretely ranked, or otherwise compact enough for the declared route action to
attain its minimum. Then there exists a unique canonical optimal route
representative
\[
 \rho_{\mathrm{opt}}^{\mathrm{can}}
 =
 \rho_{\mathrm{opt}}^{\mathrm{can}}(q_{\mathrm{clk}},q_{\mathrm{tar}}).
\]
\end{theorem}

\begin{proof}
Since the admissible route space is nonempty and the route action attains a
minimum under the stated finiteness/rank/compactness hypothesis, at least one
minimizing route exists. If several geometric representatives realize the same
minimum, the Stone--von Neumann representation identification, Teichmuller
extremal-route normalization, and QAM/MML normal-form readout identify them as
one canonical information route. The QAM/MML normal-form rule then selects one
canonical representative of that class. Hence the canonical optimal route
representative exists and is unique in the route language of the paper.
\end{proof}

\subsection{Optimal-route map}
\label{subsec:v410r10-optimal-route-map}

\begin{definition}[Clock-anchored optimal-route map]
\label{def:v410r10-clock-anchored-optimal-route-map}
For all target information loci for which the admissible route space is
nonempty, define the optimal-route map
\[
 \mathcal O_{\mathrm{clk}}:
 q_{\mathrm{tar}}
 \longmapsto
 \rho_{\mathrm{opt}}^{\mathrm{can}}
 (q_{\mathrm{clk}},q_{\mathrm{tar}}).
\]
This map is the local ``map of optimal routes'' determined by the clock anchor,
the target locus, the route action, and the canonical equivalence rule.
\end{definition}

\begin{remark}[Why this is stronger than arbitrary path choice]
\label{rem:v410r10-stronger-than-arbitrary-path-choice}
The route is not chosen by convenience. The clock anchor fixes the reference,
the target fixes the information locus, the admissible class fixes the allowed
router/tensor/guard moves, the route action fixes optimality, and the canonical
equivalence/normal-form rule removes representational multiplicity.
\end{remark}

\subsection{Effect on OP~1/4}
\label{subsec:v410r10-effect-on-op14}

\begin{definition}[Clock-to-anchor admissibility for OP~1/4]
\label{def:v410r10-clock-to-anchor-admissibility-op14}
For a comparison-relevant component \(K_i\in\mathcal K_{1/4}\), define
\[
 \mathcal R_{\mathrm{adm}}
 (q_{\mathrm{clk}},K_i\cap A_{1/4}^{\mathrm{sel}})
 \neq\varnothing
\]
to mean that at least one route-admissible clock-to-anchor route reaches a
common anchor inside \(K_i\).
\end{definition}

\begin{proposition}[OP~1/4 cover criterion becomes route existence]
\label{prop:v410r10-op14-cover-becomes-route-existence}
If for every \(K_i\in\mathcal K_{1/4}\) one has
\[
 \mathcal R_{\mathrm{adm}}
 (q_{\mathrm{clk}},K_i\cap A_{1/4}^{\mathrm{sel}})
 \neq\varnothing,
\]
then each component has a unique canonical optimal clock-to-anchor route
\[
 \rho_i^{\mathrm{opt,can}}
 =
 \rho_{\mathrm{opt}}^{\mathrm{can}}
 (q_{\mathrm{clk}},K_i\cap A_{1/4}^{\mathrm{sel}}),
\]
and the finite cover task is reduced to proving admissible route existence for
each relevant component.
\end{proposition}

\begin{proof}
Apply Theorem~\ref{thm:v410r10-clock-anchored-canonical-optimal-route-principle}
to each nonempty admissible route space
\(\mathcal R_{\mathrm{adm}}(q_{\mathrm{clk}},K_i\cap A_{1/4}^{\mathrm{sel}})\).
Each such route space has a unique canonical optimal representative. Therefore
the component-cover problem is no longer an arbitrary anchor search; it is a
finite family of clock-to-anchor admissible-route existence problems.
\end{proof}

\subsection{Effect on OP~4 and OP~5}
\label{subsec:v410r10-effect-on-op4-op5}

\begin{remark}[OP~4 as optimal-route-map stability]
\label{rem:v410r10-op4-as-optimal-route-map-stability}
With the clock-anchored optimal-route map in place, OP~4 can be read as a
stability problem for \(\mathcal O_{\mathrm{clk}}\): the question is whether the
canonical optimal routes remain transport-stable under the admitted router and
tensor transitions.
\end{remark}

\begin{proposition}[Boundary route principle for OP~5(B)]
\label{prop:v410r10-boundary-route-principle-op5b}
If the guard-preserving route space
\[
 \mathcal R_{\mathrm{adm}}^{G_{\partial}}
 (q_{\mathrm{clk}},Y_{5}^{(B),\mathrm{prom}})
\]
is nonempty, then there is a unique canonical guard-preserving boundary
promotion route
\[
 \rho_{5}^{(B),\mathrm{opt,can}}.
\]
The promotion package \(\Pi_{5}^{(B)}\) may then be read as the execution of
this canonical route, subject to the remaining HC/MM compatibility and residual
checks.
\end{proposition}

\begin{proof}
This is the guarded instance of the clock-anchored canonical optimal route
principle, with admissibility restricted by \(G_{\partial}\). The existence of
at least one guard-preserving route gives a minimizing canonical class and a
QAM/MML representative. The execution of that representative gives the
route-level reading of \(\Pi_{5}^{(B)}\); HC/MM compatibility and residual
control remain separate obligations.
\end{proof}

\subsection{reduced obstruction}
\label{subsec:v410r10-reduced-obstruction}

\begin{theorem}[OP-facing reduction]
\label{thm:v410r10-op-facing-reduction}
After the clock-anchored canonical optimal route principle, the next OP-facing
obstruction is not selection among many equally optimal routes. It is the
existence of admissible routes to the required anchor or boundary targets:
\[
 \mathcal R_{\mathrm{adm}}(q_{\mathrm{clk}},q_{\mathrm{tar}})\neq\varnothing.
\]
Once such a route exists, its canonical optimal representative is fixed by the
Stone--von Neumann/Teichmuller/QAM normal-form reading.
\end{theorem}

\begin{proof}
The preceding theorem gives a unique canonical optimal representative whenever
the admissible route space is nonempty and the declared action attains its
minimum. Therefore the ambiguity of multiple equal routes is removed at the
canonical-route level. What remains is to prove that the relevant admissible
route spaces are nonempty and that their guards and residuals are respected.
\end{proof}

\begin{remark}[OP status after]
\label{rem:v410r10-op-status}
OP~1/4, OP~4, and OP~5 remain open. The gain is that their next route-level
form is sharper: show admissible clock-to-target route existence, then compare
the current route to the canonical optimal-route map.
\end{remark}



\clearpage
\section{Admissible Route Certificate and Optimal-Route Trigger}
\label{sec:v410r11-admissible-route-certificate-optimal-route-trigger}

\begin{quote}\small
\noindent\textbf{Status.}
The layer says that once an admissible clock-to-target route exists, the
Stone--von Neumann/Teichmuller/QAM-normalized optimal route is canonically
fixed. The present layer therefore moves the OP-facing task one step
earlier: instead of searching for the optimal route directly, it defines the
finite certificate needed to prove that an admissible route exists.
\end{quote}

\subsection{Route certificate}
\label{subsec:v410r11-route-certificate}

\begin{definition}[Admissible route certificate]
\label{def:v410r11-admissible-route-certificate}
Let \(q_{\mathrm{clk}}\) be the clock anchor and \(q_{\mathrm{tar}}\) a target
information locus. An admissible route certificate is a finite record
\[
 \mathsf{Cert}_{\rho}(q_{\mathrm{clk}},q_{\mathrm{tar}})
 =
 (V_{\rho},E_{\rho},s,t,\tau,\Gamma,\Theta)
\]
where:
\begin{enumerate}[leftmargin=*]
 \item \(V_{\rho}\) is a finite family of route states;
 \item \(E_{\rho}\subseteq V_{\rho}\times V_{\rho}\) is a finite family of
 directed admissible route steps;
 \item \(s=q_{\mathrm{clk}}\) is the clock start state;
 \item \(t=q_{\mathrm{tar}}\) is the target state or target component;
 \item \(\tau\) assigns the QAM/MML/MMA tensor type to each edge;
 \item \(\Gamma\) records the active guard obligations for each edge;
 \item \(\Theta\) records the route action/rank data used by the local pass.
\end{enumerate}
\end{definition}

\begin{definition}[Certificate validity]
\label{def:v410r11-certificate-validity}
The certificate \(\mathsf{Cert}_{\rho}\) is valid if:
\begin{enumerate}[leftmargin=*]
 \item there is a directed path in \((V_{\rho},E_{\rho})\) from \(s\) to \(t\);
 \item every edge on the path has matching tensor interfaces under \(\tau\);
 \item every edge preserves the guards declared by \(\Gamma\);
 \item the route action/rank record \(\Theta\) is defined along the path.
\end{enumerate}
\end{definition}

\begin{lemma}[Certificate implies admissible route existence]
\label{lem:v410r11-certificate-implies-route-existence}
If \(\mathsf{Cert}_{\rho}(q_{\mathrm{clk}},q_{\mathrm{tar}})\) is valid, then
\[
 \mathcal R_{\mathrm{adm}}(q_{\mathrm{clk}},q_{\mathrm{tar}})
 \neq\varnothing.
\]
\end{lemma}

\begin{proof}
A valid certificate contains a directed path from the clock start state to the
target state. Each edge in the path is admissible by tensor-interface matching,
guard preservation, and route-action definability. Concatenating the certified
edges yields an admissible route from \(q_{\mathrm{clk}}\) to \(q_{\mathrm{tar}}\).
\end{proof}

\subsection{Optimal-route trigger}
\label{subsec:v410r11-optimal-route-trigger}

\begin{theorem}[Certificate triggers canonical optimal route]
\label{thm:v410r11-certificate-triggers-canonical-optimal-route}
If a valid admissible route certificate exists for
\((q_{\mathrm{clk}},q_{\mathrm{tar}})\), then the principle supplies the
canonical optimal route representative
\[
 \rho_{\mathrm{opt}}^{\mathrm{can}}
 (q_{\mathrm{clk}},q_{\mathrm{tar}}).
\]
\end{theorem}

\begin{proof}
By Lemma~\ref{lem:v410r11-certificate-implies-route-existence}, a valid
certificate proves that the admissible route space is nonempty. Applying the
clock-anchored canonical optimal route principle gives the unique canonical
optimal route representative, assuming the declared action/rank setting is one
in which the minimum is attained.
\end{proof}

\subsection{OP~1/4 certificate form}
\label{subsec:v410r11-op14-certificate-form}

\begin{definition}[OP~1/4 component certificate]
\label{def:v410r11-op14-component-certificate}
For a comparison-relevant component \(K_i\in\mathcal K_{1/4}\), an OP~1/4
component certificate is a valid route certificate for the target
\[
 q_{\mathrm{tar},i}
 :=
 K_i\cap A_{1/4}^{\mathrm{sel}}.
\]
It proves that the route system can reach a common anchor in the component
from the clock anchor.
\end{definition}

\begin{proposition}[OP~1/4 finite certificate reduction]
\label{prop:v410r11-op14-finite-certificate-reduction}
If every component \(K_i\in\mathcal K_{1/4}\) admits a valid OP~1/4 component
certificate, then every component has a canonical optimal clock-to-anchor route
\[
 \rho_i^{\mathrm{opt,can}}.
\]
Consequently the cover defect \(D_{1/4}^{\mathrm{cov}}\) is reduced to the
absence or failure of such finite certificates.
\end{proposition}

\begin{proof}
Each valid component certificate gives an admissible route from
\(q_{\mathrm{clk}}\) to \(K_i\cap A_{1/4}^{\mathrm{sel}}\). By the local pass
principle, each such route space has a canonical optimal representative. If a
component has no certificate, it remains in the certificate-level obstruction.
Thus the previous cover defect is represented by the finite certificate failure
set.
\end{proof}

\subsection{OP~5(B) boundary certificate form}
\label{subsec:v410r11-op5b-boundary-certificate-form}

\begin{definition}[Boundary promotion certificate]
\label{def:v410r11-boundary-promotion-certificate}
A boundary promotion certificate for OP~5(B) is a valid route certificate for
\[
 q_{\mathrm{tar}}
 =
 Y_{5}^{(B),\mathrm{prom}}
\]
inside the guard-restricted route class
\[
 \mathcal R_{\mathrm{adm}}^{G_{\partial}}
 (q_{\mathrm{clk}},Y_{5}^{(B),\mathrm{prom}}).
\]
Its guard record \(\Gamma\) must include
\[
 G_{\partial}(u)=G_{\partial}(v)
\]
for every certified boundary edge \(u\to v\).
\end{definition}

\begin{proposition}[Boundary certificate triggers promotion route]
\label{prop:v410r11-boundary-certificate-triggers-promotion-route}
A valid OP~5(B) boundary promotion certificate supplies a unique canonical
guard-preserving boundary route
\[
 \rho_{5}^{(B),\mathrm{opt,can}},
\]
and hence gives the route-level input needed to interpret
\(\Pi_{5}^{(B)}\) as execution of that canonical route, subject to HC/MM
compatibility and residual checks.
\end{proposition}

\begin{proof}
The certificate proves non-emptiness of the guard-restricted route space. The
guarded route principle then supplies a canonical optimal representative.
The route-level execution of this representative is exactly the route input
needed for the promotion package. The HC/MM and residual clauses remain
separate obligations.
\end{proof}

\subsection{Certificate obstruction}
\label{subsec:v410r11-certificate-obstruction}

\begin{definition}[Certificate obstruction set]
\label{def:v410r11-certificate-obstruction-set}
For a finite family of OP targets \(\{q_{\mathrm{tar},i}\}_{i\in I}\), define
\[
 D_{\rho}^{\mathrm{cert}}
 :=
 \{\,i\in I\mid
 \mathsf{Cert}_{\rho}(q_{\mathrm{clk}},q_{\mathrm{tar},i})
 \text{ is absent or invalid}\,\}.
\]
\end{definition}

\begin{theorem}[OP-facing reduction]
\label{thm:v410r11-op-facing-reduction}
After the local pass, the route-level OP-facing obstruction is the certificate obstruction
\[
 D_{\rho}^{\mathrm{cert}}.
\]
If \(D_{\rho}^{\mathrm{cert}}=\varnothing\), then all selected OP targets have
canonical optimal clock-to-target routes.
\end{theorem}

\begin{proof}
For each target, a valid certificate proves admissible route existence. By the local pass,
admissible route existence triggers the canonical optimal route representative.
Therefore precisely those targets without valid certificates remain obstructed.
The obstruction set is \(D_{\rho}^{\mathrm{cert}}\).
\end{proof}

\begin{remark}[OP status after]
\label{rem:v410r11-op-status}
OP~1/4, OP~4, and OP~5 remain open. The improvement is operational: the next
thing to build is no longer an abstract optimal route, but a finite typed
certificate proving that the route exists.
\end{remark}



\clearpage
\section{Route-Certificate Tree and Arborescence Criterion}
\label{sec:v410r12-route-certificate-tree-arborescence}

\begin{quote}\small
\noindent\textbf{Status.}
The layer reduced route existence to finite typed route certificates. The
present layer packages many target certificates into one clock-rooted
directed route tree. This is the first explicit spanning/arborescence reading
of the OP-facing route problem: the tree must start at the clock anchor, reach
all selected OP targets, respect tensor interfaces, preserve guards, and carry
route-action data. This does not close OP~1/4, OP~4, or OP~5. It makes the
next obstruction a finite tree-validity problem.
\end{quote}

\subsection{Clock-rooted certificate tree}
\label{subsec:v410r12-clock-rooted-certificate-tree}

\begin{definition}[Clock-rooted route-certificate tree]
\label{def:v410r12-clock-rooted-route-certificate-tree}
Let \(\mathcal Q=\{q_{\mathrm{tar},i}\}_{i\in I}\) be a finite family of
target information loci. A clock-rooted route-certificate tree is a finite
directed tree/arborescence
\[
 \mathsf T_{\rho}
 =
 (V_T,E_T,q_{\mathrm{clk}},\mathcal Q,\tau_T,\Gamma_T,\Theta_T)
\]
where:
\begin{enumerate}[leftmargin=*]
 \item \(q_{\mathrm{clk}}\in V_T\) is the root;
 \item every \(q_{\mathrm{tar},i}\in\mathcal Q\) is reached from
 \(q_{\mathrm{clk}}\) by a directed path in \(E_T\);
 \item \(\tau_T\) assigns tensor/interface type data to every edge;
 \item \(\Gamma_T\) assigns the active guard record to every edge;
 \item \(\Theta_T\) assigns route action/rank data to every edge or path.
\end{enumerate}
\end{definition}

\begin{definition}[Tree validity]
\label{def:v410r12-tree-validity}
The tree \(\mathsf T_{\rho}\) is valid if every root-to-target path extracted
from it is a valid admissible route certificate in the sense of
Definition~\ref{def:v410r11-certificate-validity}.
\end{definition}

\begin{lemma}[Tree validity gives all component certificates]
\label{lem:v410r12-tree-validity-gives-all-certificates}
If \(\mathsf T_{\rho}\) is valid, then every target \(q_{\mathrm{tar},i}\in
\mathcal Q\) has a valid admissible route certificate.
\end{lemma}

\begin{proof}
For each target \(q_{\mathrm{tar},i}\), follow the unique directed path in the
tree from the clock root to that target. By tree validity, this path satisfies
tensor-interface matching, guard preservation, and route-action definability.
Therefore it is a valid admissible route certificate for the target.
\end{proof}

\subsection{Spanning and Steiner readings}
\label{subsec:v410r12-spanning-steiner-readings}

\begin{definition}[Spanning route tree]
\label{def:v410r12-spanning-route-tree}
A valid route-certificate tree is a spanning route tree for \(\mathcal Q\) if
its required vertices are exactly the clock root and the selected targets,
together with any already declared route states.
\end{definition}

\begin{definition}[Steiner route tree]
\label{def:v410r12-steiner-route-tree}
A valid route-certificate tree is a Steiner route tree if it may introduce
additional admissible intermediate states
\[
 z_1,\ldots,z_\ell
\]
such as HC compensation points, MML router states, MMA readout states, boundary
collar states, or selector-anchor states, provided every additional state has a
declared tensor type and guard role.
\end{definition}

\begin{remark}[Why the Steiner version is natural here]
\label{rem:v410r12-why-steiner-natural}
The OP-facing system rarely needs only a direct edge from clock to target.
Intermediate states are often the meaningful physics: HC compensation,
boundary-collar stabilization, route-mask decoding, or selector comparison can
all appear as necessary admissible middle points. Thus the natural object is a
guarded tensor-labeled directed Steiner tree, with ordinary spanning trees as
the special case.
\end{remark}

\subsection{Cost and optimal tree}
\label{subsec:v410r12-cost-and-optimal-tree}

\begin{definition}[Tree action]
\label{def:v410r12-tree-action}
The action of a valid route-certificate tree is
\[
 \mathcal A_T(\mathsf T_{\rho})
 :=
 \sum_{e\in E_T} w(e)
 +
 \lambda\,\mathsf G_T
 +
 \mu\,\mathsf R_T,
\]
where \(w(e)\) is the edge route cost/rank, \(\mathsf G_T\) records guard
tension or guard-check complexity, and \(\mathsf R_T\) records residual terms
remaining on the tree. The constants \(\lambda,\mu\) are bookkeeping weights
for the current OP-facing readout.
\end{definition}

\begin{definition}[Canonical optimal route tree]
\label{def:v410r12-canonical-optimal-route-tree}
If the valid tree class
\[
 \mathcal T_{\mathrm{adm}}(q_{\mathrm{clk}},\mathcal Q)
\]
is nonempty, the canonical optimal route tree is the QAM/MML normal-form
representative
\[
 \mathsf T_{\rho}^{\mathrm{opt,can}}
 :=
 \operatorname*{argmin}_{\mathsf T_{\rho}\in
 \mathcal T_{\mathrm{adm}}(q_{\mathrm{clk}},\mathcal Q)}
 \mathcal A_T(\mathsf T_{\rho})
\]
after canonical route equivalence has identified representationally equivalent
trees.
\end{definition}

\begin{proposition}[Valid tree triggers all canonical routes]
\label{prop:v410r12-valid-tree-triggers-all-canonical-routes}
If a valid route-certificate tree exists for \((q_{\mathrm{clk}},\mathcal Q)\),
then every target in \(\mathcal Q\) has a canonical optimal clock-to-target
route. If the valid tree class is nonempty and the tree action attains its
minimum, then \(\mathsf T_{\rho}^{\mathrm{opt,can}}\) exists.
\end{proposition}

\begin{proof}
A valid tree gives a valid certificate for every target by
Lemma~\ref{lem:v410r12-tree-validity-gives-all-certificates}. By and the local pass,
each certificate triggers a canonical optimal route. If the whole valid tree
class is nonempty and the tree action attains its minimum, then minimization
over canonical tree classes produces a canonical optimal route tree.
\end{proof}

\subsection{OP~1/4, OP~4, and OP~5(B) readings}
\label{subsec:v410r12-op-readings}

\begin{definition}[OP~1/4 anchor-cover tree]
\label{def:v410r12-op14-anchor-cover-tree}
An OP~1/4 anchor-cover tree is a valid route-certificate tree with target
family
\[
 \mathcal Q_{1/4}
 :=
 \{K_i\cap A_{1/4}^{\mathrm{sel}}\mid K_i\in\mathcal K_{1/4}\}.
\]
It certifies clock-to-anchor reachability for every comparison-relevant
component.
\end{definition}

\begin{definition}[OP~5(B) boundary-promotion tree]
\label{def:v410r12-op5b-boundary-promotion-tree}
An OP~5(B) boundary-promotion tree is a valid route-certificate tree in the
guard-restricted class
\[
 \mathcal T_{\mathrm{adm}}^{G_{\partial}}
 (q_{\mathrm{clk}},Y_5^{(B),\mathrm{prom}})
\]
such that every boundary edge preserves \(G_{\partial}\) and the terminal path
reaches the promoted boundary target.
\end{definition}

\begin{remark}[OP~4 as tree stability]
\label{rem:v410r12-op4-as-tree-stability}
OP~4 is now naturally read as stability of the certified route tree or of the
clock-anchored optimal-route map extracted from it. Informally, OP~1/4 builds
or certifies the anchor tree, OP~4 tests transport stability of that tree, and
OP~5(B) tests whether a boundary-promotion tree can be built without breaking
the boundary guard.
\end{remark}

\subsection{Tree obstruction}
\label{subsec:v410r12-tree-obstruction}

\begin{definition}[Tree obstruction]
\label{def:v410r12-tree-obstruction}
For a selected target family \(\mathcal Q\), define the tree obstruction
\[
 D_T^{\mathrm{cert}}
 :=
 \begin{cases}
 \varnothing, & \text{if a valid route-certificate tree exists},\\
 \mathcal Q_{\mathrm{fail}}, & \text{otherwise, recording targets not reached
 by any valid tree path.}
 \end{cases}
\]
\end{definition}

\begin{theorem}[OP-facing tree reduction]
\label{thm:v410r12-op-facing-tree-reduction}
After the local pass, the route-level OP-facing obstruction can be represented by the
tree obstruction \(D_T^{\mathrm{cert}}\). If
\[
 D_T^{\mathrm{cert}}=\varnothing,
\]
then the selected targets are reached by a valid tensor-typed, guard-preserving
clock-rooted route tree and hence possess canonical optimal route
representatives.
\end{theorem}

\begin{proof}
If \(D_T^{\mathrm{cert}}=\varnothing\), a valid route-certificate tree exists.
By Proposition~\ref{prop:v410r12-valid-tree-triggers-all-canonical-routes}, all
targets then have valid certificates and canonical optimal routes. Conversely,
any target not reached by a valid tree path remains a certificate failure and is
recorded in the tree obstruction.
\end{proof}

\begin{remark}[OP status after]
\label{rem:v410r12-op-status}
OP~1/4, OP~4, and OP~5(B) remain open. The gain is that their route-level
form is now a finite tree/arborescence problem with tensor types, guards, and
route action visible.
\end{remark}



\clearpage
\section{Routing-Protocol Analogies for OP-Facing Trees}
\label{sec:v410r13-routing-protocol-analogies}

\begin{quote}\small
\noindent\textbf{Status.}
This layer extracts routing ideas from established Internet and network-routing
protocol families. The extraction is architectural only: OSPF, IS-IS, BGP,
MPLS/Traffic Engineering, Segment Routing, RSVP-TE, RPL, BFD, PCEP, and
controller-based routing are not imported as physics and are not proof
certificates. They supply robust design patterns for the QAM/MML route-tree
certificate layer introduced in the local proof chain.
\end{quote}

\subsection{Protocol-pattern table}
\label{subsec:v410r13-protocol-pattern-table}

\begin{center}
\small
\setlength{\tabcolsep}{3pt}
\begin{tabular}{@{}p{0.18\textwidth}P{0.34\textwidth}P{0.38\textwidth}@{}}
\toprule
Protocol idea & Network role & OP-facing route-tree analogue\\
\midrule
OSPF / IS-IS & Link-state map and shortest-path computation inside a domain &
Build a local map of admissible tensor/guard edges and compute a clock-rooted
route tree.\\
BGP & Path-vector routing with policy constraints between domains &
Record whole route histories and policy/guard filters before accepting a
route.\\
MPLS / TE & Traffic-engineered paths with explicit constraints &
Treat HC/MML/MMA/collar constraints as traffic-engineering constraints on
allowed OP paths.\\
Segment Routing & Source/route instruction list through selected waypoints &
Encode an OP route as a finite segment list of required compensator/router/readout
states.\\
RSVP-TE & Path reservation and resource/constraint signaling &
Reserve guard budget, boundary collar, or residual tolerance along a route.\\
RPL & Directed acyclic graph rooted at a sink/root in constrained networks &
Use the clock anchor as root and build a directed acyclic route tree to targets.\\
BFD / keepalive & Fast failure detection for links/routes &
Detect broken guard/tensor links before the OP certificate is accepted.\\
PCEP / SDN & Controller computes constrained paths from global view &
Let the QAM/MML router compute or audit the route tree from declared global
certificate data.\\
\bottomrule
\end{tabular}
\end{center}

\begin{remark}[Not a protocol import]
\label{rem:v410r13-not-protocol-import}
The table is a design dictionary, not a dependency. It says which proof-route
pattern resembles a known network-routing pattern. It does not assert that a
physical system uses Internet protocols.
\end{remark}

\subsection{Link-state route database}
\label{subsec:v410r13-link-state-route-database}

\begin{definition}[OP link-state database]
\label{def:v410r13-op-link-state-database}
An OP link-state database is a finite record
\[
 \mathsf{LS}_{\mathrm{OP}}
 =
 (V,E,\tau,\Gamma,\omega,\mathrm{up})
\]
where \(V\) are route states, \(E\) are possible directed steps, \(\tau\) is
the tensor/interface label, \(\Gamma\) is the guard record, \(\omega\) is the
edge rank/cost, and \(\mathrm{up}(e)\in\{0,1\}\) records whether the edge is
currently admissible.
\end{definition}

\begin{definition}[Link-state admissible edge]
\label{def:v410r13-link-state-admissible-edge}
An edge \(e:u\to v\) is link-state admissible if
\[
 \mathrm{up}(e)=1,
 \qquad
 \tau(e)\ \text{matches at its interfaces},
 \qquad
 \Gamma(e)\ \text{preserves all active guards}.
\]
\end{definition}

\begin{proposition}[Link-state tree extraction]
\label{prop:v410r13-link-state-tree-extraction}
If the OP link-state database contains a directed admissible path from
\(q_{\mathrm{clk}}\) to every target in \(\mathcal Q\), then it contains a valid
route-certificate tree for \((q_{\mathrm{clk}},\mathcal Q)\).
\end{proposition}

\begin{proof}
Choose, for each target, one directed admissible path from the clock root to the
target. The union of these paths contains a directed subgraph rooted at
\(q_{\mathrm{clk}}\) that reaches all targets. Removing redundant cycles, if
any, yields a directed arborescence or a directed acyclic route certificate
subtree. Since each retained edge is tensor-compatible and guard-preserving,
the resulting tree is valid in the sense.
\end{proof}

\subsection{Path-vector guard memory}
\label{subsec:v410r13-path-vector-guard-memory}

\begin{definition}[Path-vector guard record]
\label{def:v410r13-path-vector-guard-record}
A path-vector guard record for a route
\[
 \rho=(e_1,\ldots,e_n)
\]
is the sequence
\[
 \mathsf{PV}_{\Gamma}(\rho)
 :=
 \bigl((e_1,\Gamma(e_1)),\ldots,(e_n,\Gamma(e_n))\bigr).
\]
It records not merely the current endpoint, but the guard history of the route.
\end{definition}

\begin{lemma}[Path-vector guard rejection]
\label{lem:v410r13-path-vector-guard-rejection}
If any edge in \(\mathsf{PV}_{\Gamma}(\rho)\) violates an active guard, then
\(\rho\) is not an admissible route certificate.
\end{lemma}

\begin{proof}
Certificate validity requires every edge of the route path to preserve the
declared guards. A path-vector guard record exposes every edge-level guard
state. Therefore any recorded violation invalidates the certificate.
\end{proof}

\subsection{Segment routing and OP waypoints}
\label{subsec:v410r13-segment-routing-op-waypoints}

\begin{definition}[OP segment list]
\label{def:v410r13-op-segment-list}
An OP segment list is a finite route instruction sequence
\[
 \mathsf{Seg}_{\mathrm{OP}}
 =
 (z_0,z_1,\ldots,z_k)
\]
where \(z_0=q_{\mathrm{clk}}\), \(z_k=q_{\mathrm{tar}}\), and intermediate
segments \(z_j\) may be HC compensation points, MML router points, MMA readout
points, boundary-collar points, or selector-anchor points.
\end{definition}

\begin{proposition}[Segment list gives Steiner-tree skeleton]
\label{prop:v410r13-segment-list-gives-steiner-skeleton}
A valid OP segment list determines the skeleton of a directed Steiner route
tree. The intermediate segment states are precisely the permitted Steiner
points needed to make the route tensor-compatible and guard-preserving.
\end{proposition}

\begin{proof}
The segment list fixes a finite sequence of required waypoints from clock to
target. Any valid realization of each segment-to-segment subroute gives a
piece of the total route. The union of these pieces is a directed route tree or
path with intermediate nodes that were not necessarily original targets. These
nodes are Steiner points by definition.
\end{proof}

\subsection{Traffic-engineering constraints}
\label{subsec:v410r13-traffic-engineering-constraints}

\begin{definition}[OP traffic-engineering constraint]
\label{def:v410r13-op-traffic-engineering-constraint}
An OP traffic-engineering constraint is any additional route constraint of the
form
\[
 \mathsf{TE}
 =
 (\mathrm{maxResidual},\mathrm{guardBudget},\mathrm{curvatureBound},
 \mathrm{HCcompat},\mathrm{boundaryCollar})
\]
declared before a route certificate is accepted.
\end{definition}

\begin{definition}[TE-valid route tree]
\label{def:v410r13-te-valid-route-tree}
A route-certificate tree is TE-valid if it is valid in the sense and also
satisfies all declared OP traffic-engineering constraints.
\end{definition}

\begin{proposition}[TE-validity sharpens route admissibility]
\label{prop:v410r13-te-validity-sharpens-route-admissibility}
Replacing validity by TE-validity narrows the admissible tree class without
changing the tree logic. In particular, any TE-valid tree is a valid
route-certificate tree, but not every valid tree is TE-valid.
\end{proposition}

\begin{proof}
TE-validity is validity plus additional constraints. Hence it implies
ordinary validity. Since the extra constraints may fail for an otherwise valid
tree, the inclusion can be strict.
\end{proof}

\subsection{Failure detection and route convergence}
\label{subsec:v410r13-failure-detection-route-convergence}

\begin{definition}[Route heartbeat]
\label{def:v410r13-route-heartbeat}
A route heartbeat is a repeated local check
\[
 \mathsf{HB}(e)
 =
 (\tau(e)\ \mathrm{matches},\ \Gamma(e)\ \mathrm{preserved},\
 \omega(e)\ \mathrm{defined})
\]
attached to an edge \(e\). If a heartbeat fails, the edge is marked
\(\mathrm{up}(e)=0\) in the OP link-state database.
\end{definition}

\begin{definition}[Route convergence]
\label{def:v410r13-route-convergence}
Route convergence is the process of recomputing a valid route-certificate tree
after one or more admissible edges become unavailable.
\end{definition}

\begin{remark}[Why failure detection matters for OP~4]
\label{rem:v410r13-failure-detection-op4}
OP~4 can be read as transport stability of the optimal-route map. Failure
detection supplies the discrete analogue: if a tensor/guard edge fails, the
route tree must either reconverge to a valid canonical tree or record a
nonempty obstruction.
\end{remark}

\subsection{Protocol-pattern reduction}
\label{subsec:v410r13-protocol-pattern-reduction}

\begin{theorem}[Routing-protocol pattern reduction]
\label{thm:v410r13-routing-protocol-pattern-reduction}
The route-tree problem may be refined into the following protocol-pattern
tasks:
\[
\begin{array}{ll}
\text{Link-state task} & \text{build }\mathsf{LS}_{\mathrm{OP}},\\
\text{Path-vector task} & \text{record guard history }\mathsf{PV}_{\Gamma},\\
\text{Segment task} & \text{choose OP waypoints }\mathsf{Seg}_{\mathrm{OP}},\\
\text{TE task} & \text{apply residual/guard/curvature constraints},\\
\text{Heartbeat task} & \text{detect failed tensor/guard edges},\\
\text{Convergence task} & \text{recompute or certify obstruction.}
\end{array}
\]
These tasks preserve the logic while making the route-tree construction
operational.
\end{theorem}

\begin{proof}
Each protocol-pattern task refines one component of the tree certificate:
link-state supplies the finite edge database, path-vector supplies route
history, segment routing supplies permitted Steiner waypoints, traffic
engineering supplies admissibility constraints, heartbeat supplies failure
detection, and convergence supplies recomputation after failure. None of these
patterns changes the meaning of a valid route-certificate tree; they only make
its construction and audit explicit.
\end{proof}

\begin{remark}[OP status after]
\label{rem:v410r13-op-status}
OP~1/4, OP~4, and OP~5(B) remain open. The gain is that the route-tree problem
now has a practical protocol grammar: map, history, waypoints, constraints,
failure detection, and convergence.
\end{remark}



\clearpage
\section{Link-State Route-Tree Extraction}
\label{sec:v410r14-link-state-route-tree-extraction}

\begin{quote}\small
\noindent\textbf{Status.}
The layer introduced routing-protocol patterns. The present layer
turns the link-state pattern into an explicit route-tree extraction rule. Given
a finite OP link-state database, the admissible subgraph is formed by keeping
only tensor-compatible, guard-preserving, active edges. The route obstruction
is then exactly the target set not reachable from the clock root in this
admissible subgraph. If all targets are reachable, canonical clock-to-target
paths can be extracted and packaged as a route-certificate tree. This still
does not close OP~1/4, OP~4, or OP~5(B); it gives the first executable
route-protocol criterion for the certificate layer.
\end{quote}

\subsection{Admissible subgraph}
\label{subsec:v410r14-admissible-subgraph}

\begin{definition}[Admissible OP subgraph]
\label{def:v410r14-admissible-op-subgraph}
Given the OP link-state database
\[
 \mathsf{LS}_{\mathrm{OP}}
 =
 (V,E,\tau,\Gamma,\omega,\mathrm{up}),
\]
define the admissible subgraph
\[
 \mathcal G_{\mathrm{adm}}
 =
 (V,E_{\mathrm{adm}})
\]
by
\[
 E_{\mathrm{adm}}
 :=
 \{e\in E\mid
 \mathrm{up}(e)=1,\ 
 \tau(e)\ \mathrm{matches},\
 \Gamma(e)\ \mathrm{preserves\ all\ active\ guards}
 \}.
\]
\end{definition}

\begin{definition}[Reachability obstruction]
\label{def:v410r14-reachability-obstruction}
For a target family \(\mathcal Q=\{q_{\mathrm{tar},i}\}_{i\in I}\), define
\[
 D_{\mathrm{reach}}
 :=
 \{q_{\mathrm{tar},i}\in\mathcal Q
 \mid
 q_{\mathrm{tar},i}\notin
 \mathrm{Reach}_{\mathcal G_{\mathrm{adm}}}(q_{\mathrm{clk}})\}.
\]
\end{definition}

\begin{lemma}[Reachability is equivalent to certificate existence]
\label{lem:v410r14-reachability-equivalent-certificate-existence}
For a target \(q_{\mathrm{tar}}\), there exists a valid link-state route
certificate from \(q_{\mathrm{clk}}\) to \(q_{\mathrm{tar}}\) if and only if
\[
 q_{\mathrm{tar}}\in
 \mathrm{Reach}_{\mathcal G_{\mathrm{adm}}}(q_{\mathrm{clk}}).
\]
\end{lemma}

\begin{proof}
If a valid certificate exists, its edges are active, tensor-compatible, and
guard-preserving. Hence its path lies in \(E_{\mathrm{adm}}\), and the target
is reachable in \(\mathcal G_{\mathrm{adm}}\). Conversely, if the target is
reachable in \(\mathcal G_{\mathrm{adm}}\), a directed path from
\(q_{\mathrm{clk}}\) to \(q_{\mathrm{tar}}\) exists using only admissible
edges. Recording its vertices, edges, tensor labels, guard records, and rank
data gives a valid route certificate.
\end{proof}

\subsection{Canonical path extraction}
\label{subsec:v410r14-canonical-path-extraction}

\begin{definition}[Canonical extracted path]
\label{def:v410r14-canonical-extracted-path}
Assume \(q_{\mathrm{tar}}\) is reachable in \(\mathcal G_{\mathrm{adm}}\).
A canonical extracted path
\[
 \rho_{\mathrm{LS}}^{\mathrm{can}}
 (q_{\mathrm{clk}},q_{\mathrm{tar}})
\]
is the QAM/MML normal-form representative of a minimum-rank path from
\(q_{\mathrm{clk}}\) to \(q_{\mathrm{tar}}\) in \(\mathcal G_{\mathrm{adm}}\)
with respect to the edge rank \(\omega\), after canonical route equivalence has
identified representationally equivalent paths.
\end{definition}

\begin{proposition}[Reachable targets have canonical extracted paths]
\label{prop:v410r14-reachable-targets-have-canonical-extracted-paths}
Every target outside \(D_{\mathrm{reach}}\) has a canonical extracted path.
\end{proposition}

\begin{proof}
If a target is outside \(D_{\mathrm{reach}}\), it is reachable in the finite
admissible subgraph. Therefore there is at least one finite admissible path.
Minimizing the declared edge-rank sum over finite paths gives a minimum-rank
path. Canonical equivalence and QAM/MML normal form identify representational
multiplicity and select the canonical representative.
\end{proof}

\subsection{Tree extraction}
\label{subsec:v410r14-tree-extraction}

\begin{definition}[Extracted link-state route tree]
\label{def:v410r14-extracted-link-state-route-tree}
If \(D_{\mathrm{reach}}=\varnothing\), define the extracted link-state route
tree
\[
 \mathsf T_{\mathrm{LS}}^{\mathrm{can}}
\]
as the canonical acyclic reduction of the union of canonical extracted paths
\[
 \bigcup_{q_{\mathrm{tar}}\in\mathcal Q}
 \rho_{\mathrm{LS}}^{\mathrm{can}}
 (q_{\mathrm{clk}},q_{\mathrm{tar}}).
\]
Its edge labels, guard records, and rank data are inherited from
\(\mathsf{LS}_{\mathrm{OP}}\).
\end{definition}

\begin{theorem}[Link-state extraction criterion]
\label{thm:v410r14-link-state-extraction-criterion}
For a finite OP link-state database and finite target family \(\mathcal Q\),
\[
 D_{\mathrm{reach}}=\varnothing
\]
if and only if an extracted valid link-state route-certificate tree exists.
\end{theorem}

\begin{proof}
If \(D_{\mathrm{reach}}=\varnothing\), every target is reachable in
\(\mathcal G_{\mathrm{adm}}\). By
Proposition~\ref{prop:v410r14-reachable-targets-have-canonical-extracted-paths},
each target has a canonical extracted path. The acyclic reduction of the union
of these paths is a finite clock-rooted route tree reaching all targets, and
all of its retained edges are admissible by construction. Conversely, if a
valid extracted tree exists, every target is reached from the clock root by a
valid admissible path; hence no target lies in \(D_{\mathrm{reach}}\).
\end{proof}

\subsection{Traffic-engineering and segment refinement}
\label{subsec:v410r14-te-segment-refinement}

\begin{definition}[TE-filtered admissible graph]
\label{def:v410r14-te-filtered-admissible-graph}
Given a traffic-engineering constraint family \(\mathsf{TE}\), define
\[
 \mathcal G_{\mathrm{adm}}^{\mathsf{TE}}
\]
by further removing from \(\mathcal G_{\mathrm{adm}}\) every edge or path
segment violating the declared residual, guard-budget, curvature, HC/MM, or
boundary-collar constraints.
\end{definition}

\begin{definition}[Segment-constrained extraction]
\label{def:v410r14-segment-constrained-extraction}
Given an OP segment list
\[
 \mathsf{Seg}_{\mathrm{OP}}=(z_0,\ldots,z_k),
\]
a segment-constrained extracted route is a concatenation of canonical extracted
paths
\[
 \rho(z_0,z_1)\ast \rho(z_1,z_2)\ast\cdots\ast\rho(z_{k-1},z_k)
\]
inside \(\mathcal G_{\mathrm{adm}}^{\mathsf{TE}}\), provided every segment is
reachable.
\end{definition}

\begin{proposition}[Segment failure localizes obstruction]
\label{prop:v410r14-segment-failure-localizes-obstruction}
If a segment-constrained route fails, the obstruction is localized to the first
segment
\[
 (z_j,z_{j+1})
\]
for which \(z_{j+1}\notin
\mathrm{Reach}_{\mathcal G_{\mathrm{adm}}^{\mathsf{TE}}}(z_j)\).
\end{proposition}

\begin{proof}
A segment-constrained route is a finite concatenation. If it fails, at least
one segment is unreachable in the TE-filtered admissible graph. Taking the
first such segment localizes the obstruction. All earlier segments are
reachable by definition of first failure.
\end{proof}

\subsection{OP-facing consequence}
\label{subsec:v410r14-op-facing-consequence}

\begin{theorem}[route-protocol criterion]
\label{thm:v410r14-route-protocol-criterion}
After the local pass, the OP-facing route-tree obstruction can be checked by the finite
criterion
\[
 D_{\mathrm{reach}}=\varnothing
\]
in the admissible link-state graph, optionally refined to the TE-filtered and
segment-constrained graph. If the criterion holds, a valid canonical route
tree exists; if it fails, the failed target or failed segment is explicitly
identified.
\end{theorem}

\begin{proof}
The unfiltered statement is the link-state extraction criterion. TE filtering
and segment constraints only restrict the admissible graph further. Therefore
the same reachability logic applies, and failures are recorded as unreachable
targets or first failed segments.
\end{proof}

\begin{remark}[OP status after]
\label{rem:v410r14-op-status}
OP~1/4, OP~4, and OP~5(B) remain open. The gain is that the route/tree
certificate layer is now checkable by a finite link-state reachability
criterion with optional traffic-engineering and segment-routing refinements.
\end{remark}



\clearpage
\section{Segment-Route Waypoint Certificates}
\label{sec:v410r15-segment-route-waypoint-certificates}

\begin{quote}\small
\noindent\textbf{Status.}
The layer gave a finite link-state reachability criterion. The present
layer adds segment-route structure: instead of asking only whether a target
is reachable, we declare which intermediate waypoints must be respected by an
OP route. This is the controlled form of the Steiner-tree idea: HC, MML, MMA,
selector, and boundary-collar states may appear as required waypoints, but only
with declared tensor types and guard roles. OP~1/4, OP~4, and OP~5(B) remain
open.
\end{quote}

\subsection{Waypoint classes}
\label{subsec:v410r15-waypoint-classes}

\begin{definition}[OP waypoint class]
\label{def:v410r15-op-waypoint-class}
An OP waypoint class is one of the route-state types
\[
 \mathsf{WP}
 \in
 \{\mathsf{Clock},\mathsf{QAM},\mathsf{MML},\mathsf{MMA},
 \mathsf{HC},\mathsf{Selector},\mathsf{BoundaryCollar},
 \mathsf{Anchor},\mathsf{Target}\}.
\]
Each waypoint used in a segment certificate must carry a tensor/interface type
and a guard role.
\end{definition}

\begin{definition}[Typed waypoint]
\label{def:v410r15-typed-waypoint}
A typed waypoint is a tuple
\[
 z=(|z|,\mathsf{WP}(z),\tau(z),\Gamma(z))
\]
where \(|z|\) is the route state, \(\mathsf{WP}(z)\) its waypoint class,
\(\tau(z)\) its tensor/interface type, and \(\Gamma(z)\) its local guard
record.
\end{definition}

\subsection{Segment-route certificate}
\label{subsec:v410r15-segment-route-certificate}

\begin{definition}[Segment-route certificate]
\label{def:v410r15-segment-route-certificate}
A segment-route certificate from \(q_{\mathrm{clk}}\) to \(q_{\mathrm{tar}}\)
is a finite waypoint list
\[
 \mathsf{Seg}_{\rho}
 =
 (z_0,z_1,\ldots,z_k)
\]
with \(z_0=q_{\mathrm{clk}}\), \(z_k=q_{\mathrm{tar}}\), together with a
link-state certificate proving that every adjacent segment
\[
 z_j\leadsto z_{j+1}
\]
is reachable inside the currently selected admissible graph
\[
 \mathcal G_{\mathrm{adm}}^{\mathsf{TE}}.
\]
\end{definition}

\begin{definition}[Segment obstruction]
\label{def:v410r15-segment-obstruction}
The segment obstruction of \(\mathsf{Seg}_{\rho}\) is
\[
 D_{\mathrm{seg}}(\mathsf{Seg}_{\rho})
 :=
 \{j\mid
 z_{j+1}\notin
 \mathrm{Reach}_{\mathcal G_{\mathrm{adm}}^{\mathsf{TE}}}(z_j)\}.
\]
\end{definition}

\begin{lemma}[Segment certificate implies route certificate]
\label{lem:v410r15-segment-certificate-implies-route-certificate}
If
\[
 D_{\mathrm{seg}}(\mathsf{Seg}_{\rho})=\varnothing,
\]
then \(\mathsf{Seg}_{\rho}\) determines a valid admissible route certificate
from \(q_{\mathrm{clk}}\) to \(q_{\mathrm{tar}}\).
\end{lemma}

\begin{proof}
If the segment obstruction is empty, every adjacent pair of waypoints is linked
by an admissible route inside the TE-filtered admissible graph. Concatenating
these certified subroutes gives a route from \(z_0=q_{\mathrm{clk}}\) to
\(z_k=q_{\mathrm{tar}}\). Since each segment is tensor-compatible,
guard-preserving, and TE-valid, the concatenation is a valid admissible route
certificate.
\end{proof}

\subsection{OP~1/4 segment shape}
\label{subsec:v410r15-op14-segment-shape}

\begin{definition}[OP~1/4 selector-anchor segment template]
\label{def:v410r15-op14-selector-anchor-segment-template}
For a comparison-relevant component \(K_i\), the OP~1/4 selector-anchor
segment template is
\[
 \mathsf{Seg}_{1/4,i}
 =
 (q_{\mathrm{clk}},
 z_{\mathrm{QAM}},
 z_{\mathrm{MML}},
 z_{\mathrm{sel}},
 K_i\cap A_{1/4}^{\mathrm{sel}}).
\]
Here \(z_{\mathrm{QAM}}\) records the shifted/QAM normal-form shift,
\(z_{\mathrm{MML}}\) records the router/intertwiner step,
\(z_{\mathrm{sel}}\) records the selector/comparison interface, and the final
target is the common anchor region inside \(K_i\).
\end{definition}

\begin{proposition}[OP~1/4 segment criterion]
\label{prop:v410r15-op14-segment-criterion}
If every OP~1/4 segment template \(\mathsf{Seg}_{1/4,i}\) has empty segment
obstruction, then every relevant component \(K_i\) has a valid clock-to-anchor
route certificate and hence a canonical optimal clock-to-anchor route.
\end{proposition}

\begin{proof}
Apply Lemma~\ref{lem:v410r15-segment-certificate-implies-route-certificate} to
each segment template. Each valid segment certificate gives a route
certificate to \(K_i\cap A_{1/4}^{\mathrm{sel}}\). By the local proof chain
principles, each such certificate triggers a canonical optimal route.
\end{proof}

\subsection{OP~5(B) segment shape}
\label{subsec:v410r15-op5b-segment-shape}

\begin{definition}[OP~5(B) boundary-promotion segment template]
\label{def:v410r15-op5b-boundary-promotion-segment-template}
The OP~5(B) boundary-promotion segment template is
\[
 \mathsf{Seg}_{5}^{(B)}
 =
 (q_{\mathrm{clk}},
 z_{\mathrm{HC}},
 z_{\partial\mathrm{col}},
 z_{\mathrm{MMA}},
 Y_5^{(B),\mathrm{prom}}).
\]
Here \(z_{\mathrm{HC}}\) records HC compensation,
\(z_{\partial\mathrm{col}}\) records the boundary-collar guard state,
\(z_{\mathrm{MMA}}\) records the matrix/readout residual evaluator, and the
terminal target is the promoted boundary-side state.
\end{definition}

\begin{definition}[Boundary segment guard]
\label{def:v410r15-boundary-segment-guard}
The OP~5(B) segment template is boundary-valid only if every segment from
\(z_{\mathrm{HC}}\) onward preserves the boundary guard:
\[
 G_{\partial}(z_j)=G_{\partial}(z_{j+1}).
\]
\end{definition}

\begin{proposition}[OP~5(B) segment criterion]
\label{prop:v410r15-op5b-segment-criterion}
If \(\mathsf{Seg}_{5}^{(B)}\) is boundary-valid and
\[
 D_{\mathrm{seg}}(\mathsf{Seg}_{5}^{(B)})=\varnothing,
\]
then OP~5(B) has a valid guard-preserving route certificate to
\(Y_5^{(B),\mathrm{prom}}\). The remaining obligations are HC/MM compatibility
and post-promotion residual control.
\end{proposition}

\begin{proof}
Empty segment obstruction gives reachability for every adjacent segment inside
the admissible graph. Boundary-validity ensures preservation of
\(G_{\partial}\) along the boundary-sensitive part of the route. Therefore the
concatenated route is a valid guard-preserving certificate to the promoted
boundary target. The statement does not discharge HC/MM compatibility or
residual control; those remain explicit.
\end{proof}

\subsection{Segment-level OP obstruction}
\label{subsec:v410r15-segment-level-op-obstruction}

\begin{definition}[Segment-level OP obstruction]
\label{def:v410r15-segment-level-op-obstruction}
For a family of segment templates \(\mathfrak S\), define
\[
 D_{\mathrm{OP}}^{\mathrm{seg}}
 :=
 \{\mathsf{Seg}\in\mathfrak S
 \mid
 D_{\mathrm{seg}}(\mathsf{Seg})\neq\varnothing
 \text{ or a declared guard segment fails}\}.
\]
\end{definition}

\begin{theorem}[segment-route reduction]
\label{thm:v410r15-segment-route-reduction}
After the local pass, the OP-facing route obstruction can be refined from target
reachability to segment-template validity:
\[
 D_{\mathrm{OP}}^{\mathrm{seg}}=\varnothing
 \quad\Longrightarrow\quad
 D_{\mathrm{reach}}=\varnothing
\]
for the selected OP targets. If \(D_{\mathrm{OP}}^{\mathrm{seg}}\neq
\varnothing\), the first failed segment localizes the obstruction.
\end{theorem}

\begin{proof}
Each segment template with empty segment obstruction yields a valid route
certificate by Lemma~\ref{lem:v410r15-segment-certificate-implies-route-certificate}.
Thus if all selected segment templates are valid, all selected targets are
reachable and \(D_{\mathrm{reach}}\) is empty. If a segment template fails,
Definition~\ref{def:v410r15-segment-obstruction} identifies the failed adjacent
waypoint pair, localizing the obstruction.
\end{proof}

\begin{remark}[OP status after]
\label{rem:v410r15-op-status}
OP~1/4, OP~4, and OP~5(B) remain open. The gain is that the Steiner route tree
now has explicit permitted waypoint templates. The next task is to instantiate
the OP~1/4 and OP~5(B) segment lists against the current link-state database and
inspect their first failed segments.
\end{remark}



\clearpage
\section{Traffic-Engineering and Guard-Budget Filters}
\label{sec:v410r16-te-guard-budget-filters}

\begin{quote}\small
\noindent\textbf{Status.}
The layer declared OP segment templates. The present layer adds the
filter that decides whether a reachable segment is actually admissible after
residual, guard-budget, curvature, HC/MM-compatibility, and boundary-collar
constraints are applied. This is a traffic-engineering analogue only: the
constraints are QAM/MML/MMA route constraints, not imported network physics.
OP~1/4, OP~4, and OP~5(B) remain open.
\end{quote}

\subsection{TE profile}
\label{subsec:v410r16-te-profile}

\begin{definition}[OP traffic-engineering profile]
\label{def:v410r16-op-traffic-engineering-profile}
An OP traffic-engineering profile is a tuple
\[
 \mathsf{TE}_{\mathrm{OP}}
 =
 (\varepsilon_R,\varepsilon_G,\varepsilon_\kappa,
 \chi_{\mathrm{HC/MM}},\chi_{\partial\mathrm{col}},\chi_{\tau})
\]
where:
\begin{enumerate}[leftmargin=*]
 \item \(\varepsilon_R\) is the maximum allowed residual increment;
 \item \(\varepsilon_G\) is the maximum allowed guard tension;
 \item \(\varepsilon_\kappa\) is the maximum allowed curvature/route-bending
 load;
 \item \(\chi_{\mathrm{HC/MM}}\in\{0,1\}\) records HC/MM compatibility;
 \item \(\chi_{\partial\mathrm{col}}\in\{0,1\}\) records boundary-collar
 admissibility;
 \item \(\chi_{\tau}\in\{0,1\}\) records tensor-interface compatibility.
\end{enumerate}
\end{definition}

\begin{definition}[TE-valid edge]
\label{def:v410r16-te-valid-edge}
Let an edge \(e\) carry diagnostics
\[
 \mathsf d(e)
 =
 (\Delta R(e),\Delta G(e),\kappa(e),
 \chi_{\mathrm{HC/MM}}(e),
 \chi_{\partial\mathrm{col}}(e),
 \chi_{\tau}(e)).
\]
The edge is \(\mathsf{TE}_{\mathrm{OP}}\)-valid if
\[
 \Delta R(e)\leq\varepsilon_R,\qquad
 \Delta G(e)\leq\varepsilon_G,\qquad
 \kappa(e)\leq\varepsilon_\kappa,
\]
and
\[
 \chi_{\mathrm{HC/MM}}(e)=
 \chi_{\partial\mathrm{col}}(e)=
 \chi_{\tau}(e)=1
\]
whenever the corresponding compatibility clause is active for the OP segment.
\end{definition}

\subsection{Filtered graph}
\label{subsec:v410r16-filtered-graph}

\begin{definition}[TE-filtered admissible graph]
\label{def:v410r16-te-filtered-admissible-graph}
Given \(\mathcal G_{\mathrm{adm}}=(V,E_{\mathrm{adm}})\), define
\[
 \mathcal G_{\mathrm{adm}}^{\mathsf{TE}}
 =
 (V,E_{\mathrm{adm}}^{\mathsf{TE}})
\]
where
\[
 E_{\mathrm{adm}}^{\mathsf{TE}}
 :=
 \{e\in E_{\mathrm{adm}}\mid e
 \text{ is }\mathsf{TE}_{\mathrm{OP}}\text{-valid}\}.
\]
\end{definition}

\begin{lemma}[TE filtering preserves validity]
\label{lem:v410r16-te-filtering-preserves-validity}
Every route certificate in \(\mathcal G_{\mathrm{adm}}^{\mathsf{TE}}\) is a
valid route certificate in \(\mathcal G_{\mathrm{adm}}\).
\end{lemma}

\begin{proof}
The TE-filtered graph is a subgraph of the admissible graph. Therefore every
edge retained by the TE filter was already active, tensor-compatible, and
guard-preserving. The extra TE clauses only restrict the admissible class
further.
\end{proof}

\subsection{Guard-budget accumulation}
\label{subsec:v410r16-guard-budget-accumulation}

\begin{definition}[Route guard budget]
\label{def:v410r16-route-guard-budget}
For a finite route
\[
 \rho=(e_1,\ldots,e_n),
\]
define its accumulated guard tension by
\[
 \mathsf G(\rho)
 :=
 \sum_{j=1}^n \Delta G(e_j).
\]
The route is guard-budget valid if
\[
 \mathsf G(\rho)\leq \varepsilon_G^{\mathrm{tot}}.
\]
\end{definition}

\begin{definition}[Route residual budget]
\label{def:v410r16-route-residual-budget}
The accumulated residual increment of \(\rho\) is
\[
 \mathsf R(\rho)
 :=
 \sum_{j=1}^n \Delta R(e_j).
\]
The route is residual-budget valid if
\[
 \mathsf R(\rho)\leq \varepsilon_R^{\mathrm{tot}}.
\]
\end{definition}

\begin{proposition}[Budget-valid segment certificate]
\label{prop:v410r16-budget-valid-segment-certificate}
If every segment in \(\mathsf{Seg}_{\rho}\) is reachable inside
\(\mathcal G_{\mathrm{adm}}^{\mathsf{TE}}\), and the concatenated route is
guard-budget and residual-budget valid, then \(\mathsf{Seg}_{\rho}\) determines
a TE-valid admissible route certificate.
\end{proposition}

\begin{proof}
Reachability inside the TE-filtered graph gives a concatenated route whose
edges satisfy local TE constraints. Guard-budget and residual-budget validity
ensure that accumulated route tension and residual remain within the declared
global bounds. Hence the resulting certificate is TE-valid.
\end{proof}

\subsection{Boundary and HC/MM filters}
\label{subsec:v410r16-boundary-hcmm-filters}

\begin{definition}[Boundary-collar filter]
\label{def:v410r16-boundary-collar-filter}
A boundary-sensitive segment is boundary-collar valid if each boundary edge
satisfies
\[
 \chi_{\partial\mathrm{col}}(e)=1,
 \qquad
 G_{\partial}(\mathrm{src}(e))=G_{\partial}(\mathrm{dst}(e)).
\]
\end{definition}

\begin{definition}[HC/MM compatibility filter]
\label{def:v410r16-hcmm-compatibility-filter}
An HC/MM-sensitive segment is HC/MM-valid if each relevant edge satisfies
\[
 \chi_{\mathrm{HC/MM}}(e)=1
\]
and its readout diagram commutes up to the declared route tolerance.
\end{definition}

\begin{proposition}[OP~5(B) TE boundary criterion]
\label{prop:v410r16-op5b-te-boundary-criterion}
If \(\mathsf{Seg}_{5}^{(B)}\) has empty segment obstruction in
\(\mathcal G_{\mathrm{adm}}^{\mathsf{TE}}\), is boundary-collar valid, is
HC/MM-valid, and satisfies total guard and residual budgets, then it supplies a
TE-valid guard-preserving boundary promotion route certificate.
\end{proposition}

\begin{proof}
Empty TE-filtered segment obstruction gives reachability of all declared
segments. Boundary-collar validity preserves \(G_{\partial}\) along the
boundary-sensitive portion. HC/MM-validity supplies the required matrix/readout
compatibility. Budget validity controls accumulated guard tension and residual.
Thus the route certificate is TE-valid and guard-preserving.
\end{proof}

\subsection{TE-level obstruction}
\label{subsec:v410r16-te-level-obstruction}

\begin{definition}[TE obstruction]
\label{def:v410r16-te-obstruction}
For a family of segment templates \(\mathfrak S\), define the TE obstruction
\[
 D_{\mathrm{OP}}^{\mathsf{TE}}
 :=
 \{\mathsf{Seg}\in\mathfrak S\mid
 D_{\mathrm{seg}}^{\mathsf{TE}}(\mathsf{Seg})\neq\varnothing
 \text{ or a budget, boundary, HC/MM, or tensor clause fails}\}.
\]
\end{definition}

\begin{theorem}[TE-filter reduction]
\label{thm:v410r16-te-filter-reduction}
After the local pass, the OP-facing segment obstruction is sharpened from reachability to
TE-validity:
\[
 D_{\mathrm{OP}}^{\mathsf{TE}}=\varnothing
 \quad\Longrightarrow\quad
 D_{\mathrm{OP}}^{\mathrm{seg}}=\varnothing
 \quad\Longrightarrow\quad
 D_{\mathrm{reach}}=\varnothing.
\]
If \(D_{\mathrm{OP}}^{\mathsf{TE}}\neq\varnothing\), the failed clause records
whether the obstruction is reachability, tensor typing, guard budget, residual
budget, curvature load, HC/MM compatibility, or boundary-collar preservation.
\end{theorem}

\begin{proof}
TE-validity is segment validity plus the declared budget and compatibility
filters. Therefore an empty TE obstruction implies empty segment obstruction.
By the local pass, empty segment obstruction implies reachability. Conversely, any TE
failure is recorded by the first failed segment or by the first failed TE clause
attached to that segment.
\end{proof}

\begin{remark}[OP status after]
\label{rem:v410r16-op-status}
OP~1/4, OP~4, and OP~5(B) remain open. The gain is that route existence is no
longer treated as enough: a route must also remain within residual, guard,
curvature, HC/MM, tensor, and boundary-collar budgets.
\end{remark}



\clearpage
\section{Why the New QAM Narrows the OPs Faster}
\label{sec:v410r17-why-new-qam-narrows-faster}

\begin{quote}\small
\noindent\textbf{Status.}
This layer explains why the current QAM/MML/MMA route grammar narrows the
OP-facing search faster than the older QAM notation. It is an explanatory
architecture layer, not an OP closure and not a physical overclaim. Its role is
to document why the local proof chain could turn OP~1/4, OP~4, and OP~5(B) from global
existence questions into finite route, certificate, segment, and TE-filter
questions.
\end{quote}

\subsection{Old obstruction shape}
\label{subsec:v410r17-old-obstruction-shape}

\begin{definition}[Unfactored OP obstruction]
\label{def:v410r17-unfactored-op-obstruction}
An unfactored OP obstruction is an obstruction described only at the level of
the target problem, without separating:
\[
 \text{route},\quad
 \text{readout},\quad
 \text{guard},\quad
 \text{residual},\quad
 \text{tensor interface},\quad
 \text{canonical normal form}.
\]
In this form one can often say that an OP is open, but one cannot yet determine
which local clause is responsible.
\end{definition}

\begin{remark}[Why the old QAM was slower]
\label{rem:v410r17-why-old-qam-was-slower}
The older QAM notation could express many structural ideas, but it did not yet
force the OP into a typed route object. Consequently selector failure,
boundary failure, tensor-interface failure, guard failure, residual growth,
route nonexistence, and representational multiplicity could appear as one large
combined obstruction.
\end{remark}

\subsection{New QAM factorization}
\label{subsec:v410r17-new-qam-factorization}

\begin{definition}[QAM route factorization]
\label{def:v410r17-qam-route-factorization}
The current QAM/MML/MMA route grammar factors an OP-facing obstruction into the
ordered layers
\[
\begin{array}{rcl}
\text{clock anchor} &:& q_{\mathrm{clk}},\\
\text{target locus} &:& q_{\mathrm{tar}},\\
\text{admissible route space} &:&
 \mathcal R_{\mathrm{adm}}(q_{\mathrm{clk}},q_{\mathrm{tar}}),\\
\text{canonical optimal route} &:&
 \rho_{\mathrm{opt}}^{\mathrm{can}},\\
\text{certificate} &:&
 \mathsf{Cert}_{\rho},\\
\text{tree} &:&
 \mathsf T_{\rho},\\
\text{link-state graph} &:&
 \mathcal G_{\mathrm{adm}},\\
\text{segment list} &:&
 \mathsf{Seg}_{\rho},\\
\text{TE profile} &:&
 \mathsf{TE}_{\mathrm{OP}}.
\end{array}
\]
\end{definition}

\begin{proposition}[Factorization narrows failure]
\label{prop:v410r17-factorization-narrows-failure}
If an OP obstruction is expressed in the current route grammar, then failure can
be localized to one of the finite obstruction sets
\[
 D_{\mathrm{reach}},
 \qquad
 D_{\mathrm{OP}}^{\mathrm{seg}},
 \qquad
 D_{\mathrm{OP}}^{\mathsf{TE}},
 \qquad
 D_{\rho}^{\mathrm{cert}},
 \qquad
 D_T^{\mathrm{cert}}.
\]
\end{proposition}

\begin{proof}
The route grammar introduces a finite sequence of checks. Reachability fails
only in \(D_{\mathrm{reach}}\). Segment validity fails only in
\(D_{\mathrm{OP}}^{\mathrm{seg}}\). TE validity fails only in
\(D_{\mathrm{OP}}^{\mathsf{TE}}\). Certificate validity fails in
\(D_{\rho}^{\mathrm{cert}}\), and tree validity fails in
\(D_T^{\mathrm{cert}}\). Thus a global OP failure is replaced by a finite
menu of typed local failures.
\end{proof}

\subsection{Role of the Hilbert shift and tensor notation}
\label{subsec:v410r17-role-hilbert-shift-tensor-notation}

\begin{remark}[Reserved zero and clock anchoring]
\label{rem:v410r17-reserved-zero-clock-anchoring}
The shifted Hilbert-shift discipline helps because the inactive state and active
route atoms are no longer confused. The clock anchor can be treated as a known
reference object, while active route states are shifted into a normal form.
This prevents the old ambiguity in which ``no active route'', ``first active
route'', and ``reference state'' could blur together.
\end{remark}

\begin{remark}[Two-sided tensor notation]
\label{rem:v410r17-two-sided-tensor-notation}
The two-sided tensor notation helps because it separates readout/frame/language
interfaces from structure/action/field interfaces. A route edge is therefore
not just an arrow; it is a typed edge with two-sided compatibility obligations.
This makes tensor mismatch visible as its own failure mode rather than hiding it
inside a general OP obstruction.
\end{remark}

\subsection{Why could move quickly}
\label{subsec:v410r17-why-r10-r16-moved-quickly}

\begin{theorem}[New-QAM narrowing principle]
\label{thm:v410r17-new-qam-narrowing-principle}
The current QAM/MML/MMA route grammar accelerates OP localization because it
turns a global existence problem into a finite ordered diagnostic chain:
\[
\begin{array}{c}
\text{target problem}\\
\Downarrow\\
\text{clock-to-target route existence}\\
\Downarrow\\
\text{canonical optimal route}\\
\Downarrow\\
\text{finite certificate}\\
\Downarrow\\
\text{route tree}\\
\Downarrow\\
\text{link-state reachability}\\
\Downarrow\\
\text{segment waypoint validity}\\
\Downarrow\\
\text{TE/guard-budget validity}.
\end{array}
\]
\end{theorem}

\begin{proof}
Each arrow is conservative and typed. The clock-to-target route separates
reference from target. The canonical optimal route removes representational
multiplicity. The finite certificate converts route existence into checkable
data. The route tree packages many certificates. Link-state reachability
turns the tree problem into finite graph reachability. Segment waypoints
localize the required Steiner states. TE filters separate residual, guard,
curvature, HC/MM, boundary-collar, and tensor-interface failures. Therefore the
search space is progressively factored into smaller local checks.
\end{proof}

\subsection{What the old QAM could and could not do}
\label{subsec:v410r17-old-qam-could-and-could-not}

\begin{proposition}[Old QAM was expressive but not yet operational]
\label{prop:v410r17-old-qam-expressive-not-operational}
The older QAM notation could describe the structural landscape, but it did not
yet provide enough typed routing machinery to force a canonical OP diagnostic
sequence.
\end{proposition}

\begin{proof}
A notation may be expressive enough to describe many structural relations
without being operational enough to determine which local admissibility clause
fails. Before the Hilbert-shift normal form, router records, two-sided tensor
interfaces, route certificates, segment templates, and TE filters were added,
the obstruction could remain semantically correct but diagnostically
unfactored. The newer grammar adds the missing operational layers.
\end{proof}

\subsection{OP-facing interpretation}
\label{subsec:v410r17-op-facing-interpretation}

\begin{remark}[OP~1/4]
\label{rem:v410r17-op14-interpretation}
For OP~1/4, the old question ``does the selector exist?'' is now split into:
clock-to-anchor reachability, QAM/MML/selector segment validity, anchor
coverage, and residual/guard/curvature budgets.
\end{remark}

\begin{remark}[OP~4]
\label{rem:v410r17-op4-interpretation}
For OP~4, the old question ``is the route globally stable?'' is now split into:
optimal-route-map stability, tree heartbeat/failure detection, reconvergence,
and tensor/guard preservation along transported tree edges.
\end{remark}

\begin{remark}[OP~5(B)]
\label{rem:v410r17-op5b-interpretation}
For OP~5(B), the old question ``can boundary promotion happen?'' is now split
into: guard-preserving route existence, HC waypoint validity, boundary-collar
validity, MMA readout validity, HC/MM compatibility, and post-promotion
residual budget.
\end{remark}

\begin{remark}[Status after]
\label{rem:v410r17-status}
This layer explains the acceleration but does not close the OPs. The next
technical step is to instantiate the TE filters against the currently
declared OP~1/4 and OP~5(B) segment templates and read off the first nonempty
obstruction set.
\end{remark}



\clearpage
\section{Mixed-Readout Edge Compatibility Layer}
\label{sec:v410r18-mixed-readout-edge-compatibility}

\begin{quote}\small
\noindent\textbf{Status.}
This layer ties the traffic-engineering flags back to the earlier guarded
mixed QAM readout and two-sided tensor notation. A route edge is no longer
accepted merely because it is reachable and guard-budget valid. It must also
admit the mixed readout required by its role. The layer is conservative: mixed
readouts are derived compatibility-filtered packages of already admitted
component readouts, not new primitive tensors and not OP-closure certificates.
\end{quote}

\subsection{Edge readout modes}
\label{subsec:v410r18-edge-readout-modes}

\begin{definition}[Edge readout mode]
\label{def:v410r18-edge-readout-mode}
An edge readout mode is a declared component readout
\[
 R_a^m
\]
attached to a route edge, where \(a\) records the readout family and \(m\)
records its mode. The current OP-facing families include
\[
 a\in
 \{\mathrm{QAM},\mathrm{MML},\mathrm{MMA},\mathrm{HC},
 \partial\mathrm{col},\mathrm{sel},\mathrm{res},\tau\}.
\]
\end{definition}

\begin{definition}[Mixed-readout selector]
\label{def:v410r18-mixed-readout-selector}
For readout modes \(R_a^m\) and \(R_b^n\), a mixed-readout selector is a finite
compatibility tensor
\[
 \chi^{ab}_{mn}\in\{0,1\}
\]
or, more generally, a finite guard-valued selector. The mixed readout is
admitted only when the selector returns the admissible value.
\end{definition}

\begin{definition}[Edge mixed readout]
\label{def:v410r18-edge-mixed-readout}
The edge mixed readout determined by \(R_a^m\), \(R_b^n\), and
\(\chi^{ab}_{mn}\) is the derived package
\[
 \operatorname{Mix}_{\chi}^{ab}(R_a,R_b)(x)
 \quad\leadsto\quad
 \chi^{ab}_{mn}R_a^m(x)R_b^n(x).
\]
This expression is read with the two-sided contraction discipline: the selector
\(\chi\) performs the compatibility test before the paired readout is admitted.
\end{definition}

\subsection{From flags to mixed readouts}
\label{subsec:v410r18-from-r16-flags-to-mixed-readouts}

\begin{definition}[Flag-to-readout realization]
\label{def:v410r18-flag-to-readout-realization}
The compatibility flags are realized as mixed-readout selectors:
\[
 \chi_{\mathrm{HC/MM}}(e)
 :=
 \chi^{\mathrm{HC},\mathrm{MML/MMA}}_{mn}
 R_{\mathrm{HC}}^m(e)R_{\mathrm{MML/MMA}}^n(e),
\]
\[
 \chi_{\partial\mathrm{col}}(e)
 :=
 \chi^{\mathrm{HC},\partial}_{mn}
 R_{\mathrm{HC}}^m(e)R_{\partial\mathrm{col}}^n(e),
\]
\[
 \chi_{\tau}(e)
 :=
 \chi^{\mathrm{front},\mathrm{back}}_{mn}
 R_{\mathrm{front}}^m(e)R_{\mathrm{back}}^n(e).
\]
A flag is therefore no longer an uninterpreted Boolean; it is the value of a
declared mixed-readout compatibility test.
\end{definition}

\begin{remark}[Why this matters]
\label{rem:v410r18-why-this-matters}
The filter could reject a route because \(\chi_{\mathrm{HC/MM}}\),
\(\chi_{\partial\mathrm{col}}\), or \(\chi_{\tau}\) failed. The present layer
explains what those failures mean: the required pair of readout channels does
not admit a guarded mixed readout at that edge.
\end{remark}

\subsection{Mixed-readout-valid edge}
\label{subsec:v410r18-mixed-readout-valid-edge}

\begin{definition}[Required mixed readout set]
\label{def:v410r18-required-mixed-readout-set}
For an edge \(e\), let
\[
 \mathfrak M(e)
 =
 \{(a,b;m,n,\chi^{ab})\}
\]
be the finite family of mixed-readout pairs required by the edge role. Examples
include HC/MML, HC/MMA, HC/boundary-collar, selector/QAM, selector/MML, and
MMA/residual pairs.
\end{definition}

\begin{definition}[Mixed-readout-valid edge]
\label{def:v410r18-mixed-readout-valid-edge}
An edge \(e\) is mixed-readout-valid if every required pair in
\(\mathfrak M(e)\) is admitted:
\[
 \forall (a,b;m,n,\chi^{ab})\in\mathfrak M(e):
 \qquad
 \chi^{ab}_{mn}R_a^m(e)R_b^n(e)\ \mathrm{is\ admissible}.
\]
\end{definition}

\begin{lemma}[Mixed-readout validity refines TE validity]
\label{lem:v410r18-mixed-readout-validity-refines-te-validity}
If an edge is mixed-readout-valid for every compatibility flag required by its
TE profile, then the TE compatibility flags on that edge are realized by
declared readout data rather than by free Boolean assumptions.
\end{lemma}

\begin{proof}
By Definition~\ref{def:v410r18-flag-to-readout-realization}, each TE
compatibility flag is the value of a declared mixed-readout selector. If the
edge is mixed-readout-valid, every required selector admits the corresponding
readout pair. Hence the flag is supported by readout data.
\end{proof}

\subsection{OP~1/4 mixed readouts}
\label{subsec:v410r18-op14-mixed-readouts}

\begin{definition}[OP~1/4 selector mixed-readout package]
\label{def:v410r18-op14-selector-mixed-readout-package}
For the OP~1/4 segment
\[
 q_{\mathrm{clk}}\to z_{\mathrm{QAM}}\to z_{\mathrm{MML}}\to
 z_{\mathrm{sel}}\to K_i\cap A_{1/4}^{\mathrm{sel}},
\]
the selector mixed-readout package is
\[
 \mathfrak M_{1/4}
 :=
 \{
 \operatorname{Mix}_{\chi}(R_{\mathrm{QAM}},R_{\mathrm{MML}}),
 \operatorname{Mix}_{\chi}(R_{\mathrm{MML}},R_{\mathrm{sel}}),
 \operatorname{Mix}_{\chi}(R_{\mathrm{sel}},R_{\mathrm{anc}})
 \}.
\]
\end{definition}

\begin{proposition}[OP~1/4 mixed-readout refinement]
\label{prop:v410r18-op14-mixed-readout-refinement}
If an OP~1/4 segment route is TE-valid and every edge of the route is
mixed-readout-valid for \(\mathfrak M_{1/4}\), then the route does not merely
reach the anchor component; it reaches it through a compatible QAM/MML/selector
readout chain.
\end{proposition}

\begin{proof}
TE-validity supplies reachability, guard, residual, curvature, and tensor
budget control. Mixed-readout validity supplies the required compatibility
between the QAM readout, the MML router readout, the selector readout, and the
anchor readout. Therefore the route reaches the anchor component through the
declared readout chain rather than through an untyped path.
\end{proof}

\subsection{OP~5(B) mixed readouts}
\label{subsec:v410r18-op5b-mixed-readouts}

\begin{definition}[OP~5(B) boundary mixed-readout package]
\label{def:v410r18-op5b-boundary-mixed-readout-package}
For the OP~5(B) segment
\[
 q_{\mathrm{clk}}\to z_{\mathrm{HC}}\to z_{\partial\mathrm{col}}\to
 z_{\mathrm{MMA}}\to Y_5^{(B),\mathrm{prom}},
\]
the boundary mixed-readout package is
\[
 \mathfrak M_{5}^{(B)}
 :=
 \{
 \operatorname{Mix}_{\chi}(R_{\mathrm{HC}},R_{\partial\mathrm{col}}),
 \operatorname{Mix}_{\chi}(R_{\partial\mathrm{col}},R_{\mathrm{MMA}}),
 \operatorname{Mix}_{\chi}(R_{\mathrm{MMA}},R_{\mathrm{res}}),
 \operatorname{Mix}_{\chi}(R_{\mathrm{HC}},R_{\mathrm{MML/MM}})
 \}.
\]
\end{definition}

\begin{proposition}[OP~5(B) mixed-readout refinement]
\label{prop:v410r18-op5b-mixed-readout-refinement}
If an OP~5(B) boundary segment route is TE-valid, boundary-collar valid,
HC/MM-valid, and mixed-readout-valid for \(\mathfrak M_{5}^{(B)}\), then the
route supplies a compatible boundary-promotion readout chain. The remaining
obligation is the post-promotion residual control, not the mere existence of a
compatible readout route.
\end{proposition}

\begin{proof}
TE-validity and boundary-collar validity supply the route, guard, and boundary
constraints. HC/MM-validity supplies the matrix/readout compatibility. The
mixed-readout package asserts that HC, boundary-collar, MMA, residual, and
MML/MM readouts can be paired through declared selectors. Thus the boundary
promotion route is compatible at the readout level. The residual after
promotion remains a separate quantity to be controlled.
\end{proof}

\subsection{Mixed-readout obstruction}
\label{subsec:v410r18-mixed-readout-obstruction}

\begin{definition}[Mixed-readout obstruction]
\label{def:v410r18-mixed-readout-obstruction}
For a route \(\rho\), define
\[
 D_{\mathrm{mix}}(\rho)
 :=
 \{(e,a,b;m,n)\mid
 (a,b;m,n,\chi^{ab})\in\mathfrak M(e)
 \text{ and }
 \chi^{ab}_{mn}R_a^m(e)R_b^n(e)
 \text{ is not admissible}\}.
\]
\end{definition}

\begin{theorem}[mixed-readout reduction]
\label{thm:v410r18-mixed-readout-reduction}
After the local pass, the OP-facing TE obstruction refines to the mixed-readout
obstruction:
\[
\begin{gathered}
D_{\mathrm{mix}}(\rho)=\varnothing\quad\Longrightarrow\quad
\begin{array}{c}
\text{all required compatibility flags on }\rho\text{ are realized}\\
\text{by declared readout selectors.}
\end{array}
\end{gathered}
\]
If \(D_{\mathrm{mix}}(\rho)\neq\varnothing\), the obstruction records the exact
edge and readout pair that failed.
\end{theorem}

\begin{proof}
The implication follows directly from the flag-to-readout realization. A
nonempty mixed-readout obstruction records an edge and a required readout pair
whose selector does not admit the mixed readout. Hence the route is no longer
rejected by an opaque compatibility flag; it is rejected by a named readout-pair
failure.
\end{proof}

\begin{remark}[OP status after]
\label{rem:v410r18-op-status}
OP~1/4, OP~4, and OP~5(B) remain open. The gain is that the compatibility part
of the route filter is now readable as mixed-readout admissibility. This should
make the next OP pass more precise: one can test whether the obstruction is
reachability, segment structure, budget, or a specific failed readout mixture.
\end{remark}



\clearpage
\section{Mixed-Readout Index Registry and Reader Atlas}
\label{sec:v410r19-mixed-readout-index-registry}

\begin{quote}\small
\noindent\textbf{Status.}
This layer is a reader-facing and audit-facing index registry for the local pass
mixed-readout edge compatibility layer. It does not add new physics and does
not close OP~1/4, OP~4, or OP~5(B). It fixes the meaning of the index tuple
\[
 (e,a,b;m,n)
\]
so that later OP diagnostics can say exactly which edge, readout family, and
mode pairing failed.
\end{quote}

\subsection{Four index levels}
\label{subsec:v410r19-four-index-levels}

\begin{definition}[Mixed-readout index tuple]
\label{def:v410r19-mixed-readout-index-tuple}
A mixed-readout diagnostic index is a tuple
\[
 \iota_{\mathrm{mix}}
 :=
 (e,a,b;m,n),
\]
where:
\begin{enumerate}[leftmargin=*]
 \item \(e\) is the route edge being tested;
 \item \(a,b\) are readout-family indices;
 \item \(m,n\) are mode indices inside the corresponding readout families.
\end{enumerate}
The tuple is read as: on edge \(e\), readout family \(a\) in mode \(m\) is
being mixed with readout family \(b\) in mode \(n\).
\end{definition}

\begin{definition}[Readout-family registry]
\label{def:v410r19-readout-family-registry}
The current readout-family registry is
\[
 \mathcal A_{\mathrm{RO}}
 :=
 \{
 \mathrm{QAM},\mathrm{MML},\mathrm{MMA},\mathrm{HC},
 \partial\mathrm{col},\mathrm{sel},\mathrm{anc},\mathrm{res},\tau
 \}.
\]
The intended readings are:
\begin{center}
\small
\begin{tabular}{@{}ll@{}}
\toprule
Family & Reading\\
\midrule
\(\mathrm{QAM}\) & object, language, normal-form, and store-side readout\\
\(\mathrm{MML}\) & route, instruction, and router/intertwiner readout\\
\(\mathrm{MMA}\) & matrix, execution, and readout-semantics evaluator\\
\(\mathrm{HC}\) & Heisenberg-compensator/frame-normalization readout\\
\(\partial\mathrm{col}\) & boundary-collar and boundary-guard readout\\
\(\mathrm{sel}\) & selector/comparison-surface readout\\
\(\mathrm{anc}\) & anchor/common-target readout\\
\(\mathrm{res}\) & residual/parity/obstruction readout\\
\(\tau\) & tensor-interface/type readout\\
\bottomrule
\end{tabular}
\end{center}
\end{definition}

\begin{definition}[Mode index registry]
\label{def:v410r19-mode-index-registry}
For every readout family \(a\in\mathcal A_{\mathrm{RO}}\), the mode index
\(m\in M_a\) records the local mode in which that family is read. In these
mode sets remain named finite local registries rather than globally numeric
registers. This prevents premature global enumeration before the OP-facing
mode families have stabilized.
\end{definition}

\begin{remark}[Why modes are not globally numbered yet]
\label{rem:v410r19-why-modes-not-global-numbered}
A global numeric mode registry would be useful later, but it would be too early
here. The current OP-facing work first needs to learn which modes are actually
used by OP~1/4, OP~4, and OP~5(B). Therefore fixes the role of the mode
indices without freezing a final numeric list.
\end{remark}

\subsection{Selector grammar}
\label{subsec:v410r19-selector-grammar}

\begin{definition}[Selector-indexed mixed readout]
\label{def:v410r19-selector-indexed-mixed-readout}
Given the diagnostic index \((e,a,b;m,n)\), the corresponding selector-indexed
mixed readout is
\[
 \operatorname{Mix}_{\chi}^{ab}[m,n](e)
 :=
 \chi^{ab}_{mn}(e)\,
 R_a^m(e)\,
 R_b^n(e).
\]
The expression is admissible only if the selector permits the pair:
\[
 \chi^{ab}_{mn}(e)=1
\]
or, in a guard-valued registry, if \(\chi^{ab}_{mn}(e)\) returns an admissible
guard state.
\end{definition}

\begin{definition}[Mixed-readout obstruction with indices]
\label{def:v410r19-mixed-readout-obstruction-with-indices}
For a route \(\rho\), the indexed mixed-readout obstruction is
\[
 D_{\mathrm{mix}}(\rho)
 :=
 \{(e,a,b;m,n)\mid
 \operatorname{Mix}_{\chi}^{ab}[m,n](e)
 \text{ is not admissible}\}.
\]
\end{definition}

\begin{proposition}[Indexed failure localization]
\label{prop:v410r19-indexed-failure-localization}
If \((e,a,b;m,n)\in D_{\mathrm{mix}}(\rho)\), then the mixed-readout failure is
localized to one route edge, two readout families, and one pair of local modes.
\end{proposition}

\begin{proof}
This is immediate from the definition of \(D_{\mathrm{mix}}(\rho)\). The
tuple records exactly the tested edge, the two readout families, and the two
modes. Hence the failure is not a global compatibility failure but a local
selector-indexed readout-pair failure.
\end{proof}

\subsection{OP~1/4 atlas}
\label{subsec:v410r19-op14-atlas}

\begin{definition}[OP~1/4 mixed-readout index atlas]
\label{def:v410r19-op14-mixed-readout-index-atlas}
For the OP~1/4 segment
\[
 q_{\mathrm{clk}}\to z_{\mathrm{QAM}}\to z_{\mathrm{MML}}\to
 z_{\mathrm{sel}}\to K_i\cap A_{1/4}^{\mathrm{sel}},
\]
the default mixed-readout index atlas is
\[
\begin{array}{ccl}
e_{QM} &:& (\mathrm{QAM},\mathrm{MML};m_Q,m_M),\\
e_{MS} &:& (\mathrm{MML},\mathrm{sel};m_M,m_S),\\
e_{SA} &:& (\mathrm{sel},\mathrm{anc};m_S,m_A).
\end{array}
\]
Thus the OP~1/4 route is read as a chain of three selector-mediated tests:
QAM/MML, MML/selector, and selector/anchor.
\end{definition}

\begin{remark}[Human reading of the OP~1/4 atlas]
\label{rem:v410r19-human-reading-op14-atlas}
The OP~1/4 atlas says: first the QAM normal-form readout must be compatible
with the MML router readout; then the router readout must be compatible with
the selector readout; then the selector readout must be compatible with the
anchor readout. A failure in any of the three pairs tells which part of the
selector route must be repaired.
\end{remark}

\subsection{OP~5(B) atlas}
\label{subsec:v410r19-op5b-atlas}

\begin{definition}[OP~5(B) mixed-readout index atlas]
\label{def:v410r19-op5b-mixed-readout-index-atlas}
For the OP~5(B) segment
\[
 q_{\mathrm{clk}}\to z_{\mathrm{HC}}\to z_{\partial\mathrm{col}}\to
 z_{\mathrm{MMA}}\to Y_5^{(B),\mathrm{prom}},
\]
the default mixed-readout index atlas is
\[
\begin{array}{ccl}
e_{H\partial} &:& (\mathrm{HC},\partial\mathrm{col};m_H,m_\partial),\\
e_{\partial M} &:& (\partial\mathrm{col},\mathrm{MMA};m_\partial,m_M),\\
e_{MR} &:& (\mathrm{MMA},\mathrm{res};m_M,m_R),\\
e_{HMM} &:& (\mathrm{HC},\mathrm{MML/MM};m_H,m_{MM}).
\end{array}
\]
The last pair records the HC/MM compatibility readout that may be attached to
the boundary-promotion route even when it is not a separate visible edge.
\end{definition}

\begin{remark}[Human reading of the OP~5(B) atlas]
\label{rem:v410r19-human-reading-op5b-atlas}
The OP~5(B) atlas says: HC must be compatible with the boundary collar;
the boundary collar must be compatible with the MMA readout; the MMA readout
must be compatible with the residual evaluator; and HC/MM compatibility must
remain available as a route-level readout. This separates boundary failure from
matrix-readout failure and from residual-control failure.
\end{remark}

\subsection{Index-to-filter bridge}
\label{subsec:v410r19-index-to-filter-bridge}

\begin{proposition}[Index bridge to TE flags]
\label{prop:v410r19-index-bridge-to-te-flags}
The TE flags
\[
 \chi_{\mathrm{HC/MM}},\qquad
 \chi_{\partial\mathrm{col}},\qquad
 \chi_{\tau}
\]
are instances of the indexed selector grammar
\[
 \chi^{ab}_{mn}(e)
\]
after the readout families and modes have been chosen from the registry.
\end{proposition}

\begin{proof}
The layer realized the TE flags as mixed-readout selectors. The present
registry fixes the general selector grammar by naming the edge \(e\), readout
families \(a,b\), and mode pair \(m,n\). Therefore each TE flag is a special
case of \(\chi^{ab}_{mn}(e)\) with the relevant family and mode assignment.
\end{proof}

\begin{theorem}[mixed-readout index clarification]
\label{thm:v410r19-mixed-readout-index-clarification}
After the local pass, a mixed-readout compatibility failure is expressible in the canonical
form
\[
 (e,a,b;m,n)\in D_{\mathrm{mix}}(\rho).
\]
This means that the obstruction is localized to a specific route edge, a
specific pair of readout families, and a specific pair of local modes.
\end{theorem}

\begin{proof}
By Definition~\ref{def:v410r19-mixed-readout-index-tuple}, every mixed-readout
diagnostic has the form \((e,a,b;m,n)\). By
Definition~\ref{def:v410r19-mixed-readout-obstruction-with-indices}, failed
mixed-readout tests are precisely recorded in \(D_{\mathrm{mix}}(\rho)\).
Thus every compatibility failure can be written in the stated canonical form.
\end{proof}

\begin{remark}[OP status after]
\label{rem:v410r19-op-status}
OP~1/4, OP~4, and OP~5(B) remain open. The gain is explanatory and diagnostic:
the new mixed-readout layer is now indexed clearly enough for a reader or audit
script to understand exactly what is being tested on each route edge.
\end{remark}



\clearpage
\section{First Mixed-Readout Mode Registry}
\label{sec:v410r20-first-mixed-readout-mode-registry}

\begin{quote}\small
\noindent\textbf{Status.}
This layer continues the index registry by naming the first local mode sets
used by OP~1/4 and OP~5(B). The modes are intentionally not converted into a
final global numeric registry. They are local OP-facing modes that stabilize
the meaning of \(m,n\) in the diagnostic tuple
\[
 (e,a,b;m,n).
\]
This is a reader/audit layer and a preparation for sharper OP diagnostics; it
does not close OP~1/4, OP~4, or OP~5(B).
\end{quote}

\subsection{Local mode sets}
\label{subsec:v410r20-local-mode-sets}

\begin{definition}[Local mode set]
\label{def:v410r20-local-mode-set}
For a readout family \(a\), a local mode set \(M_a^{\mathrm{loc}}\) is the
finite set of named modes currently needed by the OP-facing route grammar. A
local mode is not yet a global numerical index; it is a stable semantic label
for a readout use.
\end{definition}

\begin{definition}[First OP-facing mode registry]
\label{def:v410r20-first-op-facing-mode-registry}
The first OP-facing mode registry is the finite family
\[
 \mathcal M_{\mathrm{OP}}^{(1)}
 =
 \{M_{\mathrm{QAM}},M_{\mathrm{MML}},M_{\mathrm{MMA}},
 M_{\mathrm{HC}},M_{\partial\mathrm{col}},
 M_{\mathrm{sel}},M_{\mathrm{anc}},M_{\mathrm{res}},M_{\tau}\}
\]
with the following local mode labels:
\[
\begin{array}{rcl}
M_{\mathrm{QAM}} &=&
\{\mathrm{normal},\mathrm{store},\mathrm{mask},\mathrm{front}\},\\
M_{\mathrm{MML}} &=&
\{\mathrm{router},\mathrm{intertwiner},\mathrm{segment},\mathrm{path}\},\\
M_{\mathrm{MMA}} &=&
\{\mathrm{matrix},\mathrm{execution},\mathrm{residualEval},\mathrm{back}\},\\
M_{\mathrm{HC}} &=&
\{\mathrm{comp},\mathrm{frame},\mathrm{boundaryNorm},\mathrm{mmCompat}\},\\
M_{\partial\mathrm{col}} &=&
\{\mathrm{boundary},\mathrm{collar},\mathrm{guard},\mathrm{stop}\},\\
M_{\mathrm{sel}} &=&
\{\mathrm{selector},\mathrm{comparison},\mathrm{fiber},\mathrm{equalizer}\},\\
M_{\mathrm{anc}} &=&
\{\mathrm{anchor},\mathrm{component},\mathrm{commonImage},\mathrm{target}\},\\
M_{\mathrm{res}} &=&
\{\mathrm{residual},\mathrm{post},\mathrm{budget},\mathrm{obstruction}\},\\
M_{\tau} &=&
\{\mathrm{frontType},\mathrm{backType},\mathrm{contract},\mathrm{compare}\}.
\end{array}
\]
\end{definition}

\begin{remark}[Why this is still local]
\label{rem:v410r20-why-still-local}
The registry is strong enough for OP~1/4 and OP~5(B) diagnostics but weak
enough to avoid premature global numbering. A later version may turn these
labels into a numeric or binary mask registry after the OP-facing mode use has
stabilized.
\end{remark}

\subsection{OP~1/4 instantiated modes}
\label{subsec:v410r20-op14-instantiated-modes}

\begin{definition}[OP~1/4 mode assignments]
\label{def:v410r20-op14-mode-assignments}
The first OP~1/4 mixed-readout mode assignments are
\[
\begin{array}{rcl}
e_{QM} &:&
(\mathrm{QAM},\mathrm{MML};
 \mathrm{normal},\mathrm{router}),\\
e_{MS} &:&
(\mathrm{MML},\mathrm{sel};
 \mathrm{router},\mathrm{selector}),\\
e_{SA} &:&
(\mathrm{sel},\mathrm{anc};
 \mathrm{selector},\mathrm{anchor}).
\end{array}
\]
Equivalently, the OP~1/4 mixed-readout tests are
\[
 \chi^{\mathrm{QAM},\mathrm{MML}}_{\mathrm{normal},\mathrm{router}}(e_{QM}),
\]
\[
 \chi^{\mathrm{MML},\mathrm{sel}}_{\mathrm{router},\mathrm{selector}}(e_{MS}),
\]
\[
 \chi^{\mathrm{sel},\mathrm{anc}}_{\mathrm{selector},\mathrm{anchor}}(e_{SA}).
\]
\end{definition}

\begin{proposition}[OP~1/4 first mode diagnostic]
\label{prop:v410r20-op14-first-mode-diagnostic}
If all three OP~1/4 mode assignments are admitted by their selectors, then the
OP~1/4 segment route passes the first mixed-readout compatibility layer.
If one fails, the failure is localized to the corresponding pair:
QAM/MML, MML/selector, or selector/anchor.
\end{proposition}

\begin{proof}
The obstruction records failures as tuples \((e,a,b;m,n)\). Substituting
the OP~1/4 mode assignments gives exactly the three listed tests. Therefore
empty failure set across these tests means the first mixed-readout layer passes,
while a nonempty failure identifies the failing pair.
\end{proof}

\subsection{OP~5(B) instantiated modes}
\label{subsec:v410r20-op5b-instantiated-modes}

\begin{definition}[OP~5(B) mode assignments]
\label{def:v410r20-op5b-mode-assignments}
The first OP~5(B) mixed-readout mode assignments are
\[
\begin{array}{rcl}
e_{H\partial} &:&
(\mathrm{HC},\partial\mathrm{col};
 \mathrm{boundaryNorm},\mathrm{collar}),\\
e_{\partial M} &:&
(\partial\mathrm{col},\mathrm{MMA};
 \mathrm{guard},\mathrm{matrix}),\\
e_{MR} &:&
(\mathrm{MMA},\mathrm{res};
 \mathrm{residualEval},\mathrm{post}),\\
e_{HMM} &:&
(\mathrm{HC},\mathrm{MML/MM};
 \mathrm{mmCompat},\mathrm{intertwiner}).
\end{array}
\]
Equivalently, the OP~5(B) mixed-readout tests are
\[
 \chi^{\mathrm{HC},\partial\mathrm{col}}_{\mathrm{boundaryNorm},\mathrm{collar}}(e_{H\partial}),
\]
\[
 \chi^{\partial\mathrm{col},\mathrm{MMA}}_{\mathrm{guard},\mathrm{matrix}}(e_{\partial M}),
\]
\[
 \chi^{\mathrm{MMA},\mathrm{res}}_{\mathrm{residualEval},\mathrm{post}}(e_{MR}),
\]
\[
 \chi^{\mathrm{HC},\mathrm{MML/MM}}_{\mathrm{mmCompat},\mathrm{intertwiner}}(e_{HMM}).
\]
\end{definition}

\begin{proposition}[OP~5(B) first mode diagnostic]
\label{prop:v410r20-op5b-first-mode-diagnostic}
If all four OP~5(B) mode assignments are admitted by their selectors, then the
boundary-promotion route passes the first mixed-readout compatibility layer.
If one fails, the failure is localized to HC/boundary-collar,
boundary-collar/MMA, MMA/residual, or HC/MM compatibility.
\end{proposition}

\begin{proof}
As in the OP~1/4 case, records a failed selector as
\((e,a,b;m,n)\). The four assignments instantiate those indices for the
OP~5(B) boundary route. Thus pass/fail is localized to one named readout-pair
test.
\end{proof}

\subsection{Mode diagnostic obstruction}
\label{subsec:v410r20-mode-diagnostic-obstruction}

\begin{definition}[Mode diagnostic obstruction]
\label{def:v410r20-mode-diagnostic-obstruction}
For a route \(\rho\), define the first mode diagnostic obstruction
\[
 D_{\mathrm{mode}}^{(1)}(\rho)
 :=
 \{(e,a,b;m,n)\in D_{\mathrm{mix}}(\rho)
 \mid m,n\in\mathcal M_{\mathrm{OP}}^{(1)}\}.
\]
\end{definition}

\begin{theorem}[first mode-registry reduction]
\label{thm:v410r20-first-mode-registry-reduction}
After the local pass, the first OP-facing mixed-readout compatibility problem can be tested
against the named local mode registry \(\mathcal M_{\mathrm{OP}}^{(1)}\). If
\[
 D_{\mathrm{mode}}^{(1)}(\rho)=\varnothing,
\]
then the route passes the first local mode layer. If not, the obstruction is a
specific edge, readout-family pair, and named mode pair.
\end{theorem}

\begin{proof}
The definition of \(D_{\mathrm{mode}}^{(1)}\) restricts the mixed-readout
obstruction to the first OP-facing mode registry. Therefore the empty set means
no mixed-readout selector fails at the first mode layer. A nonempty set gives
the exact tuple \((e,a,b;m,n)\), with \(m,n\) now named rather than schematic.
\end{proof}

\begin{remark}[Use in the next OP pass]
\label{rem:v410r20-use-in-next-op-pass}
The next OP pass can instantiate the tests before attempting larger
route-tree or residual arguments. If the first mixed-readout mode layer fails,
the correct repair target is local. If it passes, one can continue to
reachability, TE budgets, residual control, and OP~4 stability.
\end{remark}



\clearpage
\section{First Selector Table and Mixed-Readout Diagnostic Pass}
\label{sec:v410r21-first-selector-table-diagnostic-pass}

\begin{quote}\small
\noindent\textbf{Status.}
The layer named the first local mixed-readout mode registry. The present
layer turns those named mode pairs into the first selector table. A route
passes the first mixed-readout diagnostic pass exactly when the required
selector entries \(\chi^{ab}_{mn}(e)\) are admitted on the required edges.
This is still not OP closure; it is a finite compatibility precheck before the
larger reachability, TE-budget, residual, and OP~4 stability arguments.
\end{quote}

\subsection{First selector table}
\label{subsec:v410r21-first-selector-table}

\begin{definition}[First selector table]
\label{def:v410r21-first-selector-table}
The first selector table is the partial finite table
\[
 \mathsf{Sel}^{(1)}
 \subseteq
 E\times\mathcal A_{\mathrm{RO}}\times\mathcal A_{\mathrm{RO}}
 \times M\times M
\]
whose entries have the form
\[
 \mathsf{Sel}^{(1)}(e,a,b;m,n)
 =
 \chi^{ab}_{mn}(e)
 \in\{0,1\}.
\]
An entry value \(1\) means that the mixed readout
\[
 \operatorname{Mix}_{\chi}^{ab}[m,n](e)
\]
is admitted by the first local mode layer. An entry value \(0\) means that the
pair is rejected at this layer.
\end{definition}

\begin{definition}[Required first selector entries]
\label{def:v410r21-required-first-selector-entries}
For a route \(\rho\), the required first selector entry set is
\[
 \mathsf{Req}^{(1)}(\rho)
 :=
 \{(e,a,b;m,n)\mid
 \text{the first-mode atlas requires this mixed readout on }\rho\}.
\]
\end{definition}

\begin{definition}[First selector failure set]
\label{def:v410r21-first-selector-failure-set}
The first selector failure set is
\[
 D_{\chi}^{(1)}(\rho)
 :=
 \{(e,a,b;m,n)\in \mathsf{Req}^{(1)}(\rho)\mid
 \mathsf{Sel}^{(1)}(e,a,b;m,n)\neq 1\}.
\]
\end{definition}

\begin{lemma}[Selector table equals first mode diagnostic]
\label{lem:v410r21-selector-table-equals-first-mode-diagnostic}
For the local mode registry,
\[
 D_{\chi}^{(1)}(\rho)=D_{\mathrm{mode}}^{(1)}(\rho).
\]
up to the notational convention that \(D_{\chi}^{(1)}\) records the selector
table entry and \(D_{\mathrm{mode}}^{(1)}\) records the failed mixed-readout
tuple.
\end{lemma}

\begin{proof}
Both sets are built from the same required tuples \((e,a,b;m,n)\). The first
records tuples whose selector-table value is not admitted; the second records
tuples whose mixed readout is not admissible. Since first-layer mixed-readout
admissibility is exactly the condition \(\chi^{ab}_{mn}(e)=1\), the two
obstruction records agree.
\end{proof}

\subsection{OP~1/4 selector table}
\label{subsec:v410r21-op14-selector-table}

\begin{definition}[OP~1/4 first selector table]
\label{def:v410r21-op14-first-selector-table}
The OP~1/4 first selector table consists of the entries
\[
\begin{array}{rcl}
\mathsf{Sel}^{(1)}_{1/4}(e_{QM}) &=&
\chi^{\mathrm{QAM},\mathrm{MML}}_{\mathrm{normal},\mathrm{router}}(e_{QM}),\\
\mathsf{Sel}^{(1)}_{1/4}(e_{MS}) &=&
\chi^{\mathrm{MML},\mathrm{sel}}_{\mathrm{router},\mathrm{selector}}(e_{MS}),\\
\mathsf{Sel}^{(1)}_{1/4}(e_{SA}) &=&
\chi^{\mathrm{sel},\mathrm{anc}}_{\mathrm{selector},\mathrm{anchor}}(e_{SA}).
\end{array}
\]
\end{definition}

\begin{definition}[OP~1/4 selector pass]
\label{def:v410r21-op14-selector-pass}
The OP~1/4 first selector pass condition is
\[
 \mathsf{Pass}_{1/4}^{\chi,(1)}
 :=
 \bigl[
 \mathsf{Sel}^{(1)}_{1/4}(e_{QM})=
 \mathsf{Sel}^{(1)}_{1/4}(e_{MS})=
 \mathsf{Sel}^{(1)}_{1/4}(e_{SA})=1
 \bigr].
\]
\end{definition}

\begin{proposition}[OP~1/4 first selector diagnostic]
\label{prop:v410r21-op14-first-selector-diagnostic}
If \(\mathsf{Pass}_{1/4}^{\chi,(1)}\) holds, then the OP~1/4 route passes the
first QAM/MML/selector/anchor mixed-readout compatibility check. If it fails,
the failed selector entry identifies whether the obstruction is QAM/MML,
MML/selector, or selector/anchor.
\end{proposition}

\begin{proof}
The three entries are exactly the three mixed-readout pairings required by the
OP~1/4 atlas. If all are admitted, the first compatibility check passes.
If any entry is not admitted, the failed tuple names the failed pair.
\end{proof}

\subsection{OP~5(B) selector table}
\label{subsec:v410r21-op5b-selector-table}

\begin{definition}[OP~5(B) first selector table]
\label{def:v410r21-op5b-first-selector-table}
The OP~5(B) first selector table consists of the entries
\[
\begin{array}{rcl}
\mathsf{Sel}^{(1)}_{5B}(e_{H\partial}) &=&
\chi^{\mathrm{HC},\partial\mathrm{col}}_{\mathrm{boundaryNorm},\mathrm{collar}}(e_{H\partial}),\\
\mathsf{Sel}^{(1)}_{5B}(e_{\partial M}) &=&
\chi^{\partial\mathrm{col},\mathrm{MMA}}_{\mathrm{guard},\mathrm{matrix}}(e_{\partial M}),\\
\mathsf{Sel}^{(1)}_{5B}(e_{MR}) &=&
\chi^{\mathrm{MMA},\mathrm{res}}_{\mathrm{residualEval},\mathrm{post}}(e_{MR}),\\
\mathsf{Sel}^{(1)}_{5B}(e_{HMM}) &=&
\chi^{\mathrm{HC},\mathrm{MML/MM}}_{\mathrm{mmCompat},\mathrm{intertwiner}}(e_{HMM}).
\end{array}
\]
\end{definition}

\begin{definition}[OP~5(B) selector pass]
\label{def:v410r21-op5b-selector-pass}
The OP~5(B) first selector pass condition is
\[
 \mathsf{Pass}_{5B}^{\chi,(1)}
 :=
 \bigl[
 \mathsf{Sel}^{(1)}_{5B}(e_{H\partial})=
 \mathsf{Sel}^{(1)}_{5B}(e_{\partial M})=
 \mathsf{Sel}^{(1)}_{5B}(e_{MR})=
 \mathsf{Sel}^{(1)}_{5B}(e_{HMM})=1
 \bigr].
\]
\end{definition}

\begin{proposition}[OP~5(B) first selector diagnostic]
\label{prop:v410r21-op5b-first-selector-diagnostic}
If \(\mathsf{Pass}_{5B}^{\chi,(1)}\) holds, then the OP~5(B) boundary route
passes the first HC/boundary/MMA/residual and HC/MM mixed-readout compatibility
check. If it fails, the failed selector entry identifies whether the
obstruction is HC/boundary-collar, boundary-collar/MMA, MMA/residual, or HC/MM.
\end{proposition}

\begin{proof}
The four entries are exactly the four mixed-readout pairings required by the
OP~5(B) atlas. The pass condition says all four are admitted. Any failure
therefore names a unique first-layer mixed-readout obstruction.
\end{proof}

\subsection{Interaction with TE and route checks}
\label{subsec:v410r21-interaction-with-te-route-checks}

\begin{theorem}[First selector pass as compatibility precheck]
\label{thm:v410r21-first-selector-pass-compatibility-precheck}
The first selector pass is a necessary compatibility precheck for the local pass
TE-valid route claims. More precisely, a route cannot satisfy the local pass
compatibility clauses through declared mixed readouts unless
\[
 D_{\chi}^{(1)}(\rho)=\varnothing
\]
for all required first-layer mixed-readout entries.
\end{theorem}

\begin{proof}
The layer realizes the compatibility flags as mixed-readout selector
tests. The layer names the first OP-facing modes. Therefore, if a required
first-layer selector entry is not admitted, the corresponding compatibility flag
cannot be supported by declared readout data at that edge. Hence empty
\(D_{\chi}^{(1)}\) is necessary for TE compatibility through this readout
layer.
\end{proof}

\begin{remark}[Where this helps the OPs]
\label{rem:v410r21-where-this-helps-ops}
The first selector table is useful because it is cheaper to check than the full
route/tree/residual problem. If the table fails, the OP obstruction is already
localized to one mixed-readout pair. If the table passes, the next pass can
honestly move to reachability, TE budgets, residual control, and OP~4 stability.
\end{remark}

\begin{remark}[OP status after]
\label{rem:v410r21-op-status}
OP~1/4, OP~4, and OP~5(B) remain open. The gain is that the first mixed-readout
mode registry now has a concrete selector-table pass/fail diagnostic.
\end{remark}



\clearpage
\section{Selector-to-Route Gate}
\label{sec:v410r22-selector-to-route-gate}

\begin{quote}\small
\noindent\textbf{Status.}
The layer introduced the first selector table. The present layer turns
that table into a formal route gate. Before a route is tested for link-state
reachability, segment validity, traffic-engineering budgets, residual control,
or OP~4 stability, the required first-layer mixed-readout selectors must pass.
This is a conservative ordering rule: it prevents route work from proceeding
through an edge whose readout compatibility is already rejected.
\end{quote}

\subsection{Selector gate}
\label{subsec:v410r22-selector-gate}

\begin{definition}[First selector gate]
\label{def:v410r22-first-selector-gate}
For a route \(\rho\), the first selector gate is the predicate
\[
 \mathsf{Gate}_{\chi}^{(1)}(\rho)
 :=
 [D_{\chi}^{(1)}(\rho)=\varnothing].
\]
If \(\mathsf{Gate}_{\chi}^{(1)}(\rho)\) holds, the route is selector-admissible
at the first mixed-readout mode layer. If it fails, the route is stopped before
reachability or TE-budget evaluation.
\end{definition}

\begin{definition}[Selector-gated admissible edge]
\label{def:v410r22-selector-gated-admissible-edge}
An edge \(e\) is first-selector-gated admissible if every first-layer mixed
readout required on \(e\) has selector value \(1\):
\[
 \forall (e,a,b;m,n)\in\mathsf{Req}^{(1)}(e):
 \qquad
 \chi^{ab}_{mn}(e)=1.
\]
The corresponding gated edge set is
\[
 E_{\chi}^{(1)}
 :=
 \{e\in E\mid e\text{ is first-selector-gated admissible}\}.
\]
\end{definition}

\begin{definition}[Selector-gated graph]
\label{def:v410r22-selector-gated-graph}
Given a link-state graph \(\mathcal G=(V,E)\), define the selector-gated graph
\[
 \mathcal G_{\chi}^{(1)}
 :=
 (V,E_{\chi}^{(1)}).
\]
The later admissible and TE-filtered graphs may be formed only after this gate
has been applied:
\[
 \mathcal G_{\chi,\mathrm{adm}}^{(1)}
 \subseteq
 \mathcal G_{\chi}^{(1)}.
\]
\end{definition}

\subsection{Gate ordering}
\label{subsec:v410r22-gate-ordering}

\begin{definition}[OP diagnostic order after]
\label{def:v410r22-op-diagnostic-order-after-r22}
The diagnostic order is
\[
\begin{array}{c}
D_{\chi}^{(1)}(\rho)=\varnothing\\
\Downarrow\\
D_{\mathrm{reach}}=\varnothing\\
\Downarrow\\
D_{\mathrm{OP}}^{\mathrm{seg}}=\varnothing\\
\Downarrow\\
D_{\mathrm{OP}}^{\mathsf{TE}}=\varnothing\\
\Downarrow\\
\text{residual control and OP~4 stability tests.}
\end{array}
\]
\end{definition}

\begin{lemma}[Gate failure precedes route failure]
\label{lem:v410r22-gate-failure-precedes-route-failure}
If \(D_{\chi}^{(1)}(\rho)\neq\varnothing\), then \(\rho\) cannot be accepted as
a first-layer readout-compatible OP route, regardless of later reachability or
TE-budget data.
\end{lemma}

\begin{proof}
A nonempty selector failure set contains a required mixed-readout pair that is
not admitted. Since realizes the relevant compatibility flags through
mixed-readout selectors, a route with such a failure lacks the readout
compatibility required before later route checks. Later reachability cannot
repair a failed selector on a required edge.
\end{proof}

\begin{proposition}[Selector-gated reachability is sharper than raw reachability]
\label{prop:v410r22-selector-gated-reachability-sharper}
Reachability in \(\mathcal G_{\chi,\mathrm{adm}}^{(1)}\) is a sharper condition
than reachability in the raw link-state graph. In particular, if a target is
reachable in \(\mathcal G_{\chi,\mathrm{adm}}^{(1)}\), then it is reachable by a
route whose first-layer mixed readout requirements are all admitted.
\end{proposition}

\begin{proof}
The graph \(\mathcal G_{\chi,\mathrm{adm}}^{(1)}\) contains only edges that
first passed the selector gate and then passed the ordinary admissibility
filters. Therefore any path in it is built entirely from selector-admitted
edges. Raw reachability may use edges that are graph-theoretically present but
readout-incompatible; selector-gated reachability cannot.
\end{proof}

\subsection{OP~1/4 gate}
\label{subsec:v410r22-op14-gate}

\begin{definition}[OP~1/4 selector gate]
\label{def:v410r22-op14-selector-gate}
The OP~1/4 selector gate is
\[
 \mathsf{Gate}_{1/4}^{\chi,(1)}
 :=
 \mathsf{Pass}_{1/4}^{\chi,(1)}.
\]
Equivalently,
\[
 \mathsf{Gate}_{1/4}^{\chi,(1)}
 =
 [\chi^{\mathrm{QAM},\mathrm{MML}}_{\mathrm{normal},\mathrm{router}}(e_{QM})=
 \chi^{\mathrm{MML},\mathrm{sel}}_{\mathrm{router},\mathrm{selector}}(e_{MS})=
 \chi^{\mathrm{sel},\mathrm{anc}}_{\mathrm{selector},\mathrm{anchor}}(e_{SA})=1].
\]
\end{definition}

\begin{proposition}[OP~1/4 gate consequence]
\label{prop:v410r22-op14-gate-consequence}
If \(\mathsf{Gate}_{1/4}^{\chi,(1)}\) fails, then OP~1/4 should be repaired at
the first failed mixed-readout selector before testing anchor reachability or
cover adequacy. If it passes, the OP~1/4 route may proceed to selector-gated
reachability and segment checks.
\end{proposition}

\begin{proof}
By Lemma~\ref{lem:v410r22-gate-failure-precedes-route-failure}, a failed
selector gate blocks first-layer readout-compatible route acceptance. The
failed entry identifies whether the repair target is QAM/MML, MML/selector, or
selector/anchor. If the gate passes, no first-layer selector entry blocks the
later route checks.
\end{proof}

\subsection{OP~5(B) gate}
\label{subsec:v410r22-op5b-gate}

\begin{definition}[OP~5(B) selector gate]
\label{def:v410r22-op5b-selector-gate}
The OP~5(B) selector gate is
\[
 \mathsf{Gate}_{5B}^{\chi,(1)}
 :=
 \mathsf{Pass}_{5B}^{\chi,(1)}.
\]
Equivalently, it requires the four entries
\[
 \chi^{\mathrm{HC},\partial\mathrm{col}}_{\mathrm{boundaryNorm},\mathrm{collar}},
 \quad
 \chi^{\partial\mathrm{col},\mathrm{MMA}}_{\mathrm{guard},\mathrm{matrix}},
 \quad
 \chi^{\mathrm{MMA},\mathrm{res}}_{\mathrm{residualEval},\mathrm{post}},
 \quad
 \chi^{\mathrm{HC},\mathrm{MML/MM}}_{\mathrm{mmCompat},\mathrm{intertwiner}}
\]
to be admitted on their declared route edges.
\end{definition}

\begin{proposition}[OP~5(B) gate consequence]
\label{prop:v410r22-op5b-gate-consequence}
If \(\mathsf{Gate}_{5B}^{\chi,(1)}\) fails, then OP~5(B) should be repaired at
the failed mixed-readout selector before boundary-promotion reachability,
HC/MM residual, or post-promotion residual checks. If it passes, the OP~5(B)
route may proceed to boundary-collar, TE-budget, and residual tests.
\end{proposition}

\begin{proof}
The proof is the OP~5(B) instance of the selector-gate ordering. A failed
first-layer selector blocks readout-compatible boundary-promotion routing. A
passed gate removes this first obstruction and allows the later tests to become
meaningful.
\end{proof}

\subsection{Gate obstruction}
\label{subsec:v410r22-gate-obstruction}

\begin{definition}[Selector-gate obstruction]
\label{def:v410r22-selector-gate-obstruction}
For a family of OP routes \(\mathfrak R\), define
\[
 D_{\mathrm{gate}}^{(1)}
 :=
 \{\rho\in\mathfrak R\mid D_{\chi}^{(1)}(\rho)\neq\varnothing\}.
\]
\end{definition}

\begin{theorem}[selector-to-route reduction]
\label{thm:v410r22-selector-to-route-reduction}
After the local pass, the first OP-facing route precheck is
\[
 D_{\mathrm{gate}}^{(1)}=\varnothing.
\]
If this fails, the obstruction is selector/readout compatibility, not
reachability, TE-budget, residual control, or OP~4 stability. If it holds, the
route family can be passed to the selector-gated reachability and TE layers.
\end{theorem}

\begin{proof}
The selector-gate obstruction is exactly the family of routes with nonempty
first selector failures. By the gate-ordering lemma, those routes cannot be
accepted before later route tests. Conversely, if no route fails the selector
gate, the first mixed-readout compatibility layer is cleared and the next
diagnostic layers may be applied.
\end{proof}

\begin{remark}[OP status after]
\label{rem:v410r22-op-status}
OP~1/4, OP~4, and OP~5(B) remain open. The gain is procedural: the proof search
now has a strict first gate. Do not spend work on reachability or TE budgets
for a route whose first required mixed-readout selector has already failed.
\end{remark}



\clearpage
\section{Selector-Gated Reachability Diagnostic}
\label{sec:v410r23-selector-gated-reachability}

\begin{quote}\small
\noindent\textbf{Status.}
The layer installed the first selector gate. The present layer performs
the next diagnostic pass: reachability is checked only inside the graph whose
edges have passed the first mixed-readout selector gate and the ordinary
link-state admissibility filters. Thus a reachable route is not merely a graph
path; it is a first-layer readout-compatible graph path. OP~1/4, OP~4, and
OP~5(B) remain open.
\end{quote}

\subsection{Selector-gated admissible graph}
\label{subsec:v410r23-selector-gated-admissible-graph}

\begin{definition}[Selector-gated admissible graph]
\label{def:v410r23-selector-gated-admissible-graph}
Let
\[
 \mathcal G=(V,E)
\]
be the raw OP link-state graph. The first selector-gated admissible graph is
\[
 \mathcal G_{\chi,\mathrm{adm}}^{(1)}
 :=
 (V,E_{\chi,\mathrm{adm}}^{(1)})
\]
where
\[
 E_{\chi,\mathrm{adm}}^{(1)}
 :=
 \{e\in E\mid
 e\in E_{\chi}^{(1)},\
 \mathrm{up}(e)=1,\
 \tau(e)\ \mathrm{matches},\
 \Gamma(e)\ \mathrm{preserves\ active\ guards}
 \}.
\]
\end{definition}

\begin{definition}[Selector-gated reachability obstruction]
\label{def:v410r23-selector-gated-reachability-obstruction}
For a selected target family \(\mathcal Q\), define
\[
 D_{\mathrm{reach}}^{\chi,(1)}
 :=
 \{q_{\mathrm{tar}}\in\mathcal Q\mid
 q_{\mathrm{tar}}\notin
 \mathrm{Reach}_{\mathcal G_{\chi,\mathrm{adm}}^{(1)}}(q_{\mathrm{clk}})\}.
\]
\end{definition}

\begin{lemma}[Selector-gated reachability implies ordinary reachability]
\label{lem:v410r23-selector-gated-reachability-implies-ordinary}
If
\[
 q_{\mathrm{tar}}\in
 \mathrm{Reach}_{\mathcal G_{\chi,\mathrm{adm}}^{(1)}}(q_{\mathrm{clk}}),
\]
then \(q_{\mathrm{tar}}\) is reachable in the ordinary admissible graph
\(\mathcal G_{\mathrm{adm}}\).
\end{lemma}

\begin{proof}
The selector-gated admissible graph is a subgraph of the ordinary admissible
graph. Therefore any path in the selector-gated admissible graph is also a
path in the ordinary admissible graph.
\end{proof}

\begin{proposition}[Selector-gated route certificate]
\label{prop:v410r23-selector-gated-route-certificate}
If
\[
 q_{\mathrm{tar}}\in
 \mathrm{Reach}_{\mathcal G_{\chi,\mathrm{adm}}^{(1)}}(q_{\mathrm{clk}}),
\]
then there exists a route certificate from \(q_{\mathrm{clk}}\) to
\(q_{\mathrm{tar}}\) whose edges are first-layer selector compatible,
tensor-compatible, active, and guard-preserving.
\end{proposition}

\begin{proof}
Reachability gives a directed path in
\(\mathcal G_{\chi,\mathrm{adm}}^{(1)}\). By construction, every edge of this
path is first-layer selector compatible, active, tensor-compatible, and
guard-preserving. Recording the path data gives the stated route certificate.
\end{proof}

\subsection{OP~1/4 selector-gated reachability}
\label{subsec:v410r23-op14-selector-gated-reachability}

\begin{definition}[OP~1/4 selector-gated target family]
\label{def:v410r23-op14-selector-gated-target-family}
The OP~1/4 selector-gated target family is
\[
 \mathcal Q_{1/4}^{\chi}
 :=
 \{K_i\cap A_{1/4}^{\mathrm{sel}}\mid K_i\in\mathcal K_{1/4}\}
\]
read inside \(\mathcal G_{\chi,\mathrm{adm}}^{(1)}\).
\end{definition}

\begin{definition}[OP~1/4 selector-gated reachability obstruction]
\label{def:v410r23-op14-selector-gated-reachability-obstruction}
The OP~1/4 selector-gated reachability obstruction is
\[
 D_{\mathrm{reach},1/4}^{\chi,(1)}
 :=
 \{K_i\cap A_{1/4}^{\mathrm{sel}}\in\mathcal Q_{1/4}^{\chi}
 \mid
 K_i\cap A_{1/4}^{\mathrm{sel}}\notin
 \mathrm{Reach}_{\mathcal G_{\chi,\mathrm{adm}}^{(1)}}(q_{\mathrm{clk}})\}.
\]
\end{definition}

\begin{proposition}[OP~1/4 first selector-gated reachability pass]
\label{prop:v410r23-op14-first-selector-gated-reachability-pass}
If
\[
 \mathsf{Gate}_{1/4}^{\chi,(1)}
\]
holds and
\[
 D_{\mathrm{reach},1/4}^{\chi,(1)}=\varnothing,
\]
then every relevant OP~1/4 anchor component has a first-layer
readout-compatible clock-to-anchor route certificate.
\end{proposition}

\begin{proof}
The gate condition ensures that the required OP~1/4 selector entries are
admitted. Empty selector-gated reachability obstruction says every selected
anchor target is reachable in the selector-gated admissible graph. Applying
Proposition~\ref{prop:v410r23-selector-gated-route-certificate} to each target
gives first-layer readout-compatible route certificates.
\end{proof}

\subsection{OP~5(B) selector-gated reachability}
\label{subsec:v410r23-op5b-selector-gated-reachability}

\begin{definition}[OP~5(B) selector-gated target]
\label{def:v410r23-op5b-selector-gated-target}
The OP~5(B) selector-gated target is
\[
 q_{5B}^{\mathrm{tar}}
 :=
 Y_5^{(B),\mathrm{prom}}
\]
read inside the selector-gated boundary-sensitive graph
\[
 \mathcal G_{\chi,\mathrm{adm}}^{(1),\partial}.
\]
\end{definition}

\begin{definition}[OP~5(B) selector-gated reachability obstruction]
\label{def:v410r23-op5b-selector-gated-reachability-obstruction}
The OP~5(B) selector-gated reachability obstruction is
\[
 D_{\mathrm{reach},5B}^{\chi,(1)}
 :=
 \begin{cases}
 \varnothing,
 &Y_5^{(B),\mathrm{prom}}\in
 \mathrm{Reach}_{\mathcal G_{\chi,\mathrm{adm}}^{(1),\partial}}(q_{\mathrm{clk}}),\\
 \{Y_5^{(B),\mathrm{prom}}\},
 &\text{otherwise.}
 \end{cases}
\]
\end{definition}

\begin{proposition}[OP~5(B) first selector-gated reachability pass]
\label{prop:v410r23-op5b-first-selector-gated-reachability-pass}
If
\[
 \mathsf{Gate}_{5B}^{\chi,(1)}
\]
holds and
\[
 D_{\mathrm{reach},5B}^{\chi,(1)}=\varnothing,
\]
then there exists a first-layer readout-compatible, boundary-sensitive
clock-to-promoted-target route certificate for OP~5(B). The remaining
obligations are TE budgets, HC/MM residual compatibility, post-promotion
residual control, and OP~4 stability.
\end{proposition}

\begin{proof}
The gate condition admits the required mixed-readout selector entries for the
OP~5(B) boundary route. Empty selector-gated reachability obstruction gives a
path to \(Y_5^{(B),\mathrm{prom}}\) inside the selector-gated
boundary-sensitive graph. This path yields a first-layer readout-compatible
route certificate. The statement does not discharge later budget, residual, or
stability obligations.
\end{proof}

\subsection{Selector-gated diagnostic chain}
\label{subsec:v410r23-selector-gated-diagnostic-chain}

\begin{theorem}[selector-gated reachability reduction]
\label{thm:v410r23-selector-gated-reachability-reduction}
After the local pass, the first two OP-facing route checks are
\[
 D_{\chi}^{(1)}(\rho)=\varnothing
 \qquad\text{and}\qquad
 D_{\mathrm{reach}}^{\chi,(1)}=\varnothing.
\]
If both hold for the selected OP route family, then all selected targets have
first-layer readout-compatible admissible route certificates. If the first
fails, the repair target is a selector entry. If the second fails, the repair
target is missing selector-compatible reachability.
\end{theorem}

\begin{proof}
The first condition is exactly the selector gate. The second condition
states that the selected targets are reachable after the selector gate and
ordinary admissibility filters have been applied. By the selector-gated route
certificate proposition, those reachable targets have first-layer
readout-compatible route certificates. Failure of the first condition is a
selector failure; failure of the second condition is a reachability failure
after selector filtering.
\end{proof}

\begin{remark}[OP status after]
\label{rem:v410r23-op-status}
OP~1/4, OP~4, and OP~5(B) remain open. The gain is that reachability is now
checked in the right graph: not the raw graph and not merely the admissible
graph, but the graph already filtered by first-layer mixed-readout selectors.
\end{remark}



% ---------------------------------------------------------------------------
% shifted.1.0-r61 support macros for integrated r50--r60 delta-stack bodies.
% These macros are deliberately renderer-safe: lowercase halo and exit operators are not used in display.
% ---------------------------------------------------------------------------
